% concatenated from alphaz.tex, custom.tex
\documentclass[12pt]{article}
\usepackage[T1]{fontenc}
\usepackage{lmodern}
\usepackage{microtype}
\usepackage{amsmath,amssymb,amsthm,mathtools}
\usepackage{comment}
\usepackage{cite}
\usepackage{enumitem,array}
\usepackage{authblk}
\usepackage[colorlinks=true,
            linkcolor=blue,
            citecolor=blue,
            urlcolor=blue]{hyperref}

\usepackage{geometry}
\geometry{width=170mm, top=25mm,bottom=25mm}

\numberwithin{equation}{section}

\newtheorem{theorem}{Theorem}[section]
\newtheorem{proposition}[theorem]{Proposition}
\newtheorem{lemma}[theorem]{Lemma}
\newtheorem{corollary}[theorem]{Corollary}
\theoremstyle{definition}
\newtheorem{definition}[theorem]{Definition}
\newtheorem{example}[theorem]{Example}
\theoremstyle{remark}
\newtheorem{remark}[theorem]{Remark}


\newcommand{\one}{\mathbf 1}
\newcommand{\dd}{\,\mathrm d}
\newcommand{\scal}{\operatorname{sc}}
\newcommand{\az}{$\alpha$--$z$}

\newcolumntype{L}[1]{>{\raggedright\arraybackslash}p{#1}}



%%%%%%%%%%%%%%%%%%%%%%%%%%%%%%%%    REFERENCES    ---     COLORS & FONT SIZES         %%%%%%%%%%%%%%%%%%%%%%%%%%%%%%%%%%%%%%%%%%

%%%%%%%%  Font Sizes

%\tiny
%\scriptsize
%\footnotesize
%\small
%\normalsize
%\large
%\Large
%\LARGE
%\huge
%\Huge


%%%%%%%%%%%%%%%%%%%%%%%%%%%%%%%%%%%%%%%%%%%%%%%%%%%     COLOUR   DEFINITIONS         %%%%%%%%%%%%%%%%%%%%%%%%%%%%%%%%%%%%%%%%%%%%%%%%%%%%%%%%%%%%%%%%%%%%%

%%%%%%%%%%%%%%%%%%%%  Default colors

%black, blue, brown, cyan, darkgray, gray, green, lightgray, lime, magenta, olive, orange, pink, purple, red, teal, violet, white, yellow
%

%%%%%%%%%%%%%%%%%%%%%%%$$$%%%%%%        SHORTCUTS FOR LANDAUER              %%%%%%%%%%%%%%%%%%%%%%%%%%%%%%%%%%%%%%%%%%%%%%%%%%%%%%%%%%%%


\newcommand{\textdef}[1]{{\color{purple}\textsf{#1}}}

%%%%%%%%%%%%%%%%%%%%%%%$$$%%%%%%        SHORTCUTS FOR AQFT              %%%%%%%%%%%%%%%%%%%%%%%%%%%%%%%%%%%%%%%%%%%%%%%%%%%%%%%%%%%%

\newcommand{\vNA}{von Neumann algebra }
\newcommand{\vNsA}{von Neumann subalgebra }



%%%%%%%%%%%%%%%%%%%%%%%$$$%%%%%%        SHORTCUTS FOR MATHS              %%%%%%%%%%%%%%%%%%%%%%%%%%%%%%%%%%%%%%%%%%%%%%%%%%%%%%%%%%%%


\newcommand{\rmO}{\text{O}}
\newcommand{\rmU}{\text{U}}
\newcommand{\GL}{\text{GL}}
\newcommand{\SL}{\text{SL}}
\newcommand{\SU}{\text{SU}}
\newcommand{\SO}{\text{SO}}
\newcommand{\Sp}{\text{Sp}}


\newcommand{\overbar}[1]{\mkern 1.5mu\overline{\mkern-1.5mu#1\mkern-1.5mu}\mkern 1.5mu}


%%%%%%%%%%%%%%%%%%%%%%%%%%%


\newcommand{\df}{{\rm d}}
\newcommand{\del}{\partial}
\newcommand{\wt}[1]{\widetilde{#1}}
\newcommand{\wh}[1]{\widehat{#1}}
\newcommand{\?}{\overset{?}{=}}


\newcommand{\defeq}{\mathrel{\mathop:}=}
\newcommand{\eqdef}{=\mathrel{\mathop:}}


\DeclareMathOperator{\Sch}{Sch}
\DeclareMathOperator{\sech}{sech}
\DeclareMathOperator{\csch}{csch}
\DeclareMathOperator{\arccosh}{arccosh}
\DeclareMathOperator{\arcsech}{arcsech}
\DeclareMathOperator{\arcsinh}{arcsinh}
\DeclareMathOperator{\arccsch}{arccsch}
\DeclareMathOperator{\arctanh}{arctanh}
\DeclareMathOperator{\arccoth}{arccoth}

\DeclareMathOperator{\tr}{tr}
\DeclareMathOperator{\Tr}{Tr}
\DeclareMathOperator{\rank}{rank}
\DeclareMathOperator{\diag}{diag}
\DeclareMathOperator{\Span}{span}
\DeclareMathOperator{\spec}{spec}
\DeclareMathOperator{\ad}{ad}
\DeclareMathOperator{\Ad}{Ad}
\DeclareMathOperator{\im}{im}
\DeclareMathOperator{\sign}{sign}
\DeclareMathOperator{\supp}{supp}
\renewcommand{\Re}{\operatorname{Re}}
\renewcommand{\Im}{\operatorname{Im}}

% \newcommand{\ad}[1]{\,\text{ad}_{#1}}
% \newcommand{\Ad}[1]{\,\text{Ad}_{#1}}


\newcommand{\abs}[1]{\left\vert#1\right\vert}

\newcommand{\ohgf}[1]{\tensor[_2]{F}{_1}\lB #1 \rB}   %% Ordinary Hypergeometric fn.

\newcommand{\rmnum}[1]{\romannumeral #1}
\newcommand{\Rmnum}[1]{\expandafter\@slowromancap\romannumeral #1@}



%%%%%%%%%%%%%%%%%%%%%%%%%%%%%%%%        EQUATION SHORTCUTS          %%%%%%%%%%%%%%%%%%%%%%%%%%%%%%%%%%%%%%%%%

\newcommand{\nn}{\nonumber}
\newcommand{\dsp}{\displaystyle}


\newcommand{\be}{\begin{equation}}
\newcommand{\ee}{\end{equation}}


\newcommand{\bea}[1]{\begin{align}#1\end{align}}
\newcommand{\bes}[1]{\begin{equation}\begin{split}#1 \end{split}\end{equation}}


%%%%%%%%%%%%%%%%%%%%%%%%%%%%%%%%        BRACKETS        %%%%%%%%%%%%%%%%%%%%%%%%%%%%%%%%%%%%%%%%%%%

\newcommand{\lB}{\left [}
\newcommand{\rB}{\right ]}
\newcommand{\lb}{\left (}
\newcommand{\rb}{\right )}
\newcommand{\lbr}{\left \{}						
\newcommand{\rbr}{\right \}}					
\newcommand{\lv}{\left\vert}
\newcommand{\rv}{\right\vert}
\newcommand{\lla}{\left\langle}
\newcommand{\rra}{\right\rangle}
\newcommand{\ld}{\left.}				        %% Left  continuation of bracketed expression.
\newcommand{\rd}{\right.}				        %% Right  continuation of bracketed expression.
\newcommand{\bfrac}[2]{\lb\frac{#1}{#2}\rb}  	%% Fractions in brackets.



%%%%%%%%%%%%%%%%%%%%%%%%%%%%%%%%%%%%%%%    SYMBOL SHORTCUTS     %%%%%%%%%%%%%%%%%%%%%%%%%%%%%%%%%%%%%%%%%%%%%%%%%%%%

\newcommand{\A}{\,\forall\,}
\newcommand{\C}{\mathbb{C}}
\newcommand{\I}{\infty}
\newcommand{\M}{\mathbb{M}}
\newcommand{\N}{\mathbb{N}}
\newcommand{\Qu}{\mathbb{H}}
\newcommand{\R}{\mathbb{R}}
\newcommand{\Z}{\mathbb{Z}}

%%%%%  Wide tilde

\newcommand{\wtA}{\widetilde{A}}
\newcommand{\wtB}{\widetilde{B}}
\newcommand{\wtC}{\widetilde{C}}
\newcommand{\wtD}{\widetilde{D}}
\newcommand{\wtE}{\widetilde{E}}
\newcommand{\wtF}{\widetilde{F}}
\newcommand{\wtG}{\widetilde{G}}
\newcommand{\wtH}{\widetilde{H}}
\newcommand{\wtI}{\widetilde{I}}
\newcommand{\wtJ}{\widetilde{J}}
\newcommand{\wtK}{\widetilde{K}}
\newcommand{\wtL}{\widetilde{L}}
\newcommand{\wtM}{\widetilde{M}}
\newcommand{\wtN}{\widetilde{N}}
\newcommand{\wtO}{\widetilde{O}}
\newcommand{\wtP}{\widetilde{P}}
\newcommand{\wtQ}{\widetilde{Q}}
\newcommand{\wtR}{\widetilde{R}}
\newcommand{\wtS}{\widetilde{S}}
\newcommand{\wtT}{\widetilde{T}}
\newcommand{\wtU}{\widetilde{U}}
\newcommand{\wtV}{\widetilde{V}}
\newcommand{\wtW}{\widetilde{W}}
\newcommand{\wtX}{\widetilde{X}}
\newcommand{\wtY}{\widetilde{Y}}
\newcommand{\wtZ}{\widetilde{Z}}


%%%%%  Greek Alphabet
\renewcommand{\a}{\alpha}	%%% Redefinition
\renewcommand{\b}{\beta}		%%% Redefinition
\renewcommand{\c}{\chi}		%%% Redefinition
\renewcommand{\d}{\delta}	%%% Redefinition
\newcommand{\e}{\epsilon}
\newcommand{\ve}{\varepsilon}
\newcommand{\f}{\phi}
\newcommand{\vf}{\varphi}
\newcommand{\g}{\gamma}
\newcommand{\h}{\eta}
\newcommand{\hb}{\hbar}
\newcommand{\hs}{\hslash}
\renewcommand{\i}{\iota}		%%% Redefinition
\renewcommand{\k}{\kappa}	%%% Redefinition
\renewcommand{\l}{\lambda}	%%% Redefinition
\newcommand{\mn}{\mu\nu}
\newcommand{\q}{\theta}	
\newcommand{\vq}{\vartheta}	
\renewcommand{\r}{\rho}		%%% Redefinition
\newcommand{\vr}{\varrho}	%%% Redefinition
\newcommand{\s}{\sigma}
\newcommand{\vs}{\varsigma}
\renewcommand{\t}{\tau}		%%% Redefinition
\renewcommand{\u}{\upsilon}	%%% Redefinition
\newcommand{\w}{\omega}
\newcommand{\vw}{\varomega}
\newcommand{\y}{\psi}
\newcommand{\z}{\zeta}


\newcommand{\D}{\Delta}
\newcommand{\F}{\Phi}
\newcommand{\G}{\Gamma}
\renewcommand{\L}{\Lambda}	%%% Redefinition
\newcommand{\Q}{\Theta}
\renewcommand{\S}{\Sigma}	%%% Redefinition
\newcommand{\W}{\Omega}
\newcommand{\Y}{\Psi}



%%%%%   CALIGRAPHIC
\newcommand{\cA}{\mathcal{A}}
\newcommand{\cB}{\mathcal{B}}
\newcommand{\cC}{\mathcal{C}}
\newcommand{\cD}{\mathcal{D}}
\newcommand{\cE}{\mathcal{E}}
\newcommand{\cF}{\mathcal{F}}
\newcommand{\cG}{\mathcal{G}}
\newcommand{\cH}{\mathcal{H}}
\newcommand{\cI}{\mathcal{I}}
\newcommand{\cJ}{\mathcal{J}}
\newcommand{\cK}{\mathcal{K}}
\newcommand{\cL}{\mathcal{L}}
\newcommand{\cM}{\mathcal{M}}
\newcommand{\cN}{\mathcal{N}}
\newcommand{\cO}{\mathcal{O}}
\newcommand{\cP}{\mathcal{P}}
\newcommand{\cQ}{\mathcal{Q}}
\newcommand{\cR}{\mathcal{R}}
\newcommand{\cS}{\mathcal{S}}
\newcommand{\cT}{\mathcal{T}}
\newcommand{\cU}{\mathcal{U}}
\newcommand{\cV}{\mathcal{V}}
\newcommand{\cW}{\mathcal{W}}
\newcommand{\cX}{\mathcal{X}}
\newcommand{\cY}{\mathcal{Y}}
\newcommand{\cZ}{\mathcal{Z}}

%%%%%   BLACKBOARD
\newcommand{\bA}{\mathbb{A}}
\newcommand{\bB}{\mathbb{B}}
\newcommand{\bC}{\mathbb{C}}
\newcommand{\bD}{\mathbb{D}}
\newcommand{\bE}{\mathbb{E}}
\newcommand{\bF}{\mathbb{F}}
\newcommand{\bG}{\mathbb{G}}
\newcommand{\bH}{\mathbb{H}}
\newcommand{\bI}{\mathbb{I}}
\newcommand{\bJ}{\mathbb{J}}
\newcommand{\bK}{\mathbb{K}}
\newcommand{\bL}{\mathbb{L}}
\newcommand{\bM}{\mathbb{M}}
\newcommand{\bN}{\mathbb{N}}
\newcommand{\bO}{\mathbb{O}}
\newcommand{\bP}{\mathbb{P}}
\newcommand{\bQ}{\mathbb{Q}}
\newcommand{\bR}{\mathbb{R}}
\newcommand{\bS}{\mathbb{S}}
\newcommand{\bT}{\mathbb{T}}
\newcommand{\bU}{\mathbb{U}}
\newcommand{\bV}{\mathbb{V}}
\newcommand{\bW}{\mathbb{W}}
\newcommand{\bX}{\mathbb{X}}
\newcommand{\bY}{\mathbb{Y}}
\newcommand{\bZ}{\mathbb{Z}}

%%%%%   SCRIPT
\newcommand{\sA}{\mathscr{A}}
\newcommand{\sB}{\mathscr{B}}
\newcommand{\sC}{\mathscr{C}}
\newcommand{\sD}{\mathscr{D}}
\newcommand{\sE}{\mathscr{E}}
\newcommand{\sF}{\mathscr{F}}
\newcommand{\sG}{\mathscr{G}}
\newcommand{\sH}{\mathscr{H}}
\newcommand{\sI}{\mathscr{I}}
\newcommand{\sJ}{\mathscr{J}}
\newcommand{\sK}{\mathscr{K}}
\newcommand{\sL}{\mathscr{L}}
\newcommand{\sM}{\mathscr{M}}
\newcommand{\sN}{\mathscr{N}}
\newcommand{\sO}{\mathscr{O}}
\newcommand{\sP}{\mathscr{P}}
\newcommand{\sQ}{\mathscr{Q}}
\newcommand{\sR}{\mathscr{R}}
\newcommand{\sS}{\mathscr{S}}
\newcommand{\sT}{\mathscr{T}}
\newcommand{\sU}{\mathscr{U}}
\newcommand{\sV}{\mathscr{V}}
\newcommand{\sW}{\mathscr{W}}
\newcommand{\sX}{\mathscr{X}}
\newcommand{\sY}{\mathscr{Y}}
\newcommand{\sZ}{\mathscr{Z}}

%%%%%%     Mathfrak
\newcommand{\fA}{\mathfrak{A}}
\newcommand{\fB}{\mathfrak{B}}
\newcommand{\fC}{\mathfrak{C}}
\newcommand{\fD}{\mathfrak{D}}
\newcommand{\fE}{\mathfrak{E}}
\newcommand{\fF}{\mathfrak{F}}
\newcommand{\fG}{\mathfrak{G}}
\newcommand{\fH}{\mathfrak{H}}
\newcommand{\fI}{\mathfrak{I}}
\newcommand{\fJ}{\mathfrak{J}}
\newcommand{\fK}{\mathfrak{K}}
\newcommand{\fL}{\mathfrak{L}}
\newcommand{\fM}{\mathfrak{M}}
\newcommand{\fN}{\mathfrak{N}}
\newcommand{\fO}{\mathfrak{O}}
\newcommand{\fP}{\mathfrak{P}}
\newcommand{\fQ}{\mathfrak{Q}}
\newcommand{\fR}{\mathfrak{R}}
\newcommand{\fS}{\mathfrak{S}}
\newcommand{\fT}{\mathfrak{T}}
\newcommand{\fU}{\mathfrak{U}}
\newcommand{\fV}{\mathfrak{V}}
\newcommand{\fW}{\mathfrak{W}}
\newcommand{\fX}{\mathfrak{X}}
\newcommand{\fY}{\mathfrak{Y}}
\newcommand{\fZ}{\mathfrak{Z}}
\newcommand{\fa}{\mathfrak{a}}
\newcommand{\fb}{\mathfrak{b}}
\newcommand{\fc}{\mathfrak{c}}
\newcommand{\fd}{\mathfrak{d}}
\newcommand{\fe}{\mathfrak{e}}
\newcommand{\ff}{\mathfrak{f}}
\newcommand{\fg}{\mathfrak{g}}
\newcommand{\fh}{\mathfrak{h}}
\newcommand{\fii}{\mathfrak{i}}		%%% Out of Order
\newcommand{\fj}{\mathfrak{j}}
\newcommand{\fk}{\mathfrak{k}}
\newcommand{\fl}{\mathfrak{l}}
\newcommand{\fm}{\mathfrak{m}}
\newcommand{\fn}{\mathfrak{n}}
\newcommand{\fo}{\mathfrak{o}}
\newcommand{\fp}{\mathfrak{p}}
\newcommand{\fq}{\mathfrak{q}}
\newcommand{\fr}{\mathfrak{r}}
\newcommand{\fs}{\mathfrak{s}}
\newcommand{\ft}{\mathfrak{t}}
\newcommand{\fu}{\mathfrak{u}}
\newcommand{\fv}{\mathfrak{v}}
\newcommand{\fw}{\mathfrak{w}}
\newcommand{\fx}{\mathfrak{x}}
\newcommand{\fy}{\mathfrak{y}}
\newcommand{\fz}{\mathfrak{z}}

%%%%%%%%%%%%%%%%%%%%%%%%%%%%%%%%%     END - SYMBOL DEFINITIONS       %%%%%%%%%%%%%%%%%%%%%%%%%%%%%%%%%%%%


\renewcommand{\e}{\mathrm e}


\begin{document}

\title{Fixed-ray escort representations of sandwiched and
$\alpha$--$z$ R\'enyi divergences on von Neumann algebras}

\author[1]{Tanay Kibe \thanks{tanay.kibe@ib.edu.ar}}
\affil[1]{Instituto Balseiro, Centro At\'omico Bariloche,
S.C. de Bariloche, 8400 R\'io Negro, Argentina}

\author[2]{Pratik Roy \thanks{roy.pratik92@gmail.com}}
\affil[2]{Institute of Mathematics, University of Warsaw,
ul. Banacha 2, 02-097 Warsaw, Poland}

\date{}

\maketitle

\abstract{
We represent sandwiched and $\alpha$--$z$ R\'enyi divergences as averages
of ordinary relative entropy. The $\alpha$--$z$ R\'enyi divergence is shown to be an integral over the relative entropy of a canonical family of fixed-ray escort states along the ray $z=c\alpha$.  We prove
this representation for normal states on an arbitrary von Neumann algebra, using Haagerup non-commutative $L^p$ spaces and interpolation.
The formula holds for every \(z>0\): for \(0<\alpha<1\) when the
support of the first state is contained in that of the reference
state, and for \(\alpha>1\) whenever the divergence is finite.
When the lower-order support condition fails, we identify the exact
fixed-ray support-boundary term.
The representation yields a monotone escort
profile and a convex order potential.  We use these to reformulate
one-shot testing converses, exact sandwiched strong-converse exponents, and
work-extraction reliability as signed-area or level-crossing statements, and discuss a restricted two-parameter pair-conversion rate.
}

\newpage

\setcounter{tocdepth}{2}
\tableofcontents



\section{Introduction}
\label{sec:introduction}

Relative entropy is simultaneously a measure of distinguishability, a
thermodynamic potential, and a basic geometric quantity on the state
space.  For density matrices $\rho$ and $\sigma$, it is
\begin{equation}
  D(\rho\Vert\sigma)
  =\Tr\rho(\log\rho-\log\sigma),
\end{equation}
with the usual support convention.  Araki's relative entropy gives its
intrinsic extension to normal states on an arbitrary von Neumann algebra
\cite{Araki:1976zv}.  This extension is crucial, for example in quantum field theory, when the observable algebra is Type~III and no density matrix or trace exists
on the algebra. The sandwiched R\'enyi divergences
\cite{Muller-Lennert:2013liu,Wilde:2013bdg} and the two-parameter
$\alpha$--$z$ family \cite{Audenaert:2015npv} refine relative entropy by reweighting with respect to the reference state.  The
$\alpha$--$z$ family contains the Petz line $z=1$ and the sandwiched line
$z=\alpha$.  Its extension to normal functionals on von Neumann algebras
is naturally formulated in Haagerup noncommutative $L^p$ spaces
\cite{Kato:2023aro,Kato:2023hlj,Hiai:2024qve}. These divergences have also been used to formulate energy conditions in quantum field theory \cite{Lashkari:2018nsl,Moosa:2020jwt,Roy:2022yzm,Kibe:2026wsg,Kibe:2026bcn}.

Our starting point is the refined sandwiched divergence of
Bao, Moosa, and Shehzad \cite{Bao:2019aol}.  In finite dimensions,
it was shown in \cite{Bao:2019aol} to equal the relative entropy
of a sandwiched escort state.  Integrating this differential identity
shows that the sandwiched divergence itself is an average of ordinary
relative entropies.  The same calculation suggests the correct
two-parameter extension: 
\begin{equation}
  D_{\alpha,z}(\rho\Vert\sigma)
  =\frac{\alpha}{\alpha-1}
   \int_1^\alpha
   \frac{D(\rho_{\beta}^{(c)}\Vert\sigma)}{\beta^2}\,\dd\beta,
  \qquad c=\frac z\alpha,
  \label{eq:intro-fd-formula}
\end{equation}
where $\rho_{\beta}^{(c)}$, introduced below, denotes the fixed-ray escort state 
on a characteristic ray
\begin{equation}
  z=c\alpha,\qquad c>0.
  \label{eq:intro-ray}
\end{equation}
For $0<\alpha<1$ the oriented integral is a positive average over
$[\alpha,1]$.  At $c=1$, equivalently $z=\alpha$, this specializes
exactly to the sandwiched divergence. This representation is different from the integral
representations of quantum divergences developed from hockey-stick or
layer-cake decompositions \cite{Hirche:2023caq,Liu:2025npj}.
Those formulas integrate over a likelihood threshold and reconstruct a
divergence from binary-testing data.  Here the integration variable is the
R\'enyi order itself, and the integrand is the Araki relative entropy of a
canonical escort state.  The two approaches therefore have different
integrands, geometries, and domains of application.

The principal result of this paper is that
\eqref{eq:intro-fd-formula} remains valid for normal states on an
arbitrary von Neumann algebra.  For normal states $\psi,\omega$, the integral formula requires
$s(\psi)\leq s(\omega)$.\footnote{We sometimes say that \(\psi\) is \emph{support-nested} in
\(\omega\), or simply that the pair \((\psi,\omega)\) has nested
supports, when
\(s(\psi)\leq s(\omega) \). 
Accordingly, ``nonnested'' means that this support
inclusion fails.} 
For $\alpha>1$, finiteness of the
$\alpha$--$z$ divergence already enforces this nestedness condition.  
Subject to these hypotheses, the
formula holds for every $\alpha,z>0$, $\alpha\ne1$, including parameters
outside the data-processing region.  For $0<\alpha<1$ and nonnested
supports, the escort average result is modified by an
$L^c$ support-boundary term. On the sandwiched line this reduces to the
support mass $\psi(s(\omega))$. 

The finite-dimensional differentiation argument does not extend
directly to general von Neumann algebras, where the corresponding
intermediate operators
are generally unbounded and live in order-dependent Haagerup
\(L^p\)-spaces. Our proof instead anchors the fixed ray
at one Haagerup element, transports it through a compatible weighted
$L^p$ scale, and derives an exact quotient identity.  A lower-order
$\alpha$--$z$ limit identifies the left derivative with Araki relative
entropy.  Quasi-Banach interpolation and endpoint norm continuity then
justify the fundamental theorem of calculus.  The part of this argument
with exponents below one uses the all-$p$ interpolation theorem of
Gu, Yin, and Zhang \cite{Gu:2019jfc}.

Our integral representation is useful because it turns an order-dependent
R\'enyi quantity into a probability average of ordinary relative
entropies.  Intrinsically, the escort relative entropy is nondecreasing
along every fixed ray and is the slope of a convex potential.  Operational
variational formulas consequently become signed areas, with the optimizing
order determined by an escort level crossing.  We state the implications of our representation for exact strong-converse and
cutoff-rate formulas under the hypotheses of the corresponding operational
theorems, and for work and pair-conversion applications
in the finite-dimensional settings in which they are established.

The rest of this paper is organized as follows.  Section~\ref{sec:finite-dimensional}
derives the matrix formulas.  Section~\ref{sec:preliminaries} introduces
the Haagerup $L^p$-space framework and the algebraic divergences.  All general theorem
statements and their consequences appear in
Section~\ref{sec:results-and-consequences}.  The technical proofs are deferred: Section~\ref{sec:proof-lower-srd} gives a direct proof of
the lower-order\footnote{We generally refer to the range $\alpha\in(0,1)$ as the lower order range of either sandwiched or $\alpha$--$z$ divergence, and the range $\alpha>1$ as the upper order range.} sandwiched formula in its data-processing range, and
Section~\ref{sec:proof-alpha-z} proves the full fixed-ray theorem.  Thus the
latter proof also supplies an alternative sandwiched proof and extends its
analytic identity to $0<\alpha<1/2$. Finally, we summarize our results and discuss future directions in Section~\ref{sec:discussion}.

\section{Finite-dimensional integral representations}
\label{sec:finite-dimensional}

Throughout this section, $\cH$ is a finite dimensional Hilbert space, and $\rho$ and
$\sigma$ are density matrices on $\cH$.  We first take $\sigma>0$ and,
when differentiating, $\rho>0$.  Singular states follow by approximation
and support compression, as explained after the calculation.

\subsection{The fixed-ray \texorpdfstring{$\alpha$--$z$}{alpha-z}
representation}

For $\alpha,z>0$, $\alpha\ne1$, define the {\az} moment and divergence respectively as
\begin{align}
  Q_{\alpha,z}(\rho\|\sigma)
  &\defeq
  \Tr\!\left(
    \sigma^{\frac{1-\alpha}{2z}}
    \rho^{\frac\alpha z}
    \sigma^{\frac{1-\alpha}{2z}}
  \right)^z,
  \label{eq:fd-alpha-z-moment}\\
  D_{\alpha,z}(\rho\Vert\sigma)
  &\defeq \frac{1}{\alpha-1}
  \log Q_{\alpha,z}(\rho\|\sigma).
  \label{eq:fd-alpha-z-divergence}
\end{align}
Consider a fixed ray in $(\alpha,z)$-parameter space with slope $c>0$, i.e., the ray $(\beta,c\beta)$.  For $\beta>0$, set
\begin{align}
  X_{\beta}^{(c)}
  &\defeq
  \sigma^{\frac{1-\beta}{2c\beta}}
  \rho^{1/c}
  \sigma^{\frac{1-\beta}{2c\beta}},
  \label{eq:fd-fixed-ray-factor}\\
  Q_{\beta}^{(c)}(\rho\|\sigma)
  &\defeq
  \Tr\!\left[(X_{\beta}^{(c)})^{c\beta}\right]
  =Q_{\beta,c\beta}(\rho\|\sigma),
  \label{eq:fd-fixed-ray-moment}\\
  \rho_{\beta}^{(c)}
  &\defeq
  \frac{(X_{\beta}^{(c)})^{c\beta}}
       {Q_{\beta}^{(c)}(\rho\|\sigma)}.
  \label{eq:fd-fixed-ray-escort}
\end{align}
The state $\rho_{\beta}^{(c)}$ will be called the \emph{fixed-ray
escort state}, and the family
$\beta\mapsto\rho_{\beta}^{(c)}$ the \emph{fixed-ray escort trajectory}.
Notice that $Q_1^{(c)}=1$ and
$\rho_1^{(c)}=\rho$.

Put $m_c(\beta)=Q_{\beta}^{(c)}(\rho\|\sigma)$ and
$g_c(\beta)=\log m_c(\beta)$.  Functional calculus gives
\begin{equation}
  \frac{\dd X_{\beta}^{(c)}}{\dd\beta}
  =-\frac{1}{2c\beta^2}
  \left[(\log\sigma)X_{\beta}^{(c)}
       +X_{\beta}^{(c)}(\log\sigma)\right].
  \label{eq:fd-factor-derivative}
\end{equation}
Using tracial differentiation and cyclicity,
\begin{equation}
  m_c'(\beta)
  =c\Tr\!\left[(X_{\beta}^{(c)})^{c\beta}
                  \log X_{\beta}^{(c)}\right]
   -\frac1\beta\Tr\!\left[(X_{\beta}^{(c)})^{c\beta}
                  \log\sigma\right].
  \label{eq:fd-moment-derivative}
\end{equation}
On the other hand,
\begin{equation}
  \log\rho_{\beta}^{(c)}
  =c\beta\log X_{\beta}^{(c)}-g_c(\beta),
\end{equation}
and therefore
\begin{equation}
  D(\rho_{\beta}^{(c)}\Vert\sigma)
  =\beta g_c'(\beta)-g_c(\beta).
  \label{eq:fd-refined-identity}
\end{equation}
Equivalently,
\begin{equation}
  \beta^2\frac{\dd}{\dd\beta}
  \left[
    \frac{\beta-1}{\beta}
    D_{\beta,c\beta}(\rho\Vert\sigma)
  \right]
  =D(\rho_{\beta}^{(c)}\Vert\sigma).
  \label{eq:fd-ray-refined-identity}
\end{equation}
Since the quantity in square brackets vanishes at $\beta=1$,
integration gives the desired formula.


We now extend the preceding notation to singular density matrices.  Put
$P=s(\sigma)$.  Positive powers of $\sigma$ are understood by the usual
functional calculus, with $0^u=0$ for $u>0$.  Whenever a negative power
of $\sigma$ occurs, it denotes the inverse power of the faithful
restriction $\sigma_P\defeq\sigma\vert_{P\cH}$ on $P\cH$, extended by
zero on $P^\perp\cH$.  Thus, for $0<\alpha<1$, the expression
\eqref{eq:fd-alpha-z-moment} is defined for arbitrary $\rho$ and
$\sigma$, with the convention that
$D_{\alpha,z}(\rho\Vert\sigma)=+\infty$ if
$Q_{\alpha,z}(\rho\|\sigma)=0$.  For $\alpha>1$ the generalized-inverse reading of $\sigma^{(1-\alpha)/2z}$ is used only
when $s(\rho)\le s(\sigma)$; without support inclusion it would return a finite value of
the wrong sign, or $-\infty$. We therefore set
\[
  Q_{\alpha,z}(\rho\|\sigma)
  =D_{\alpha,z}(\rho\Vert\sigma)
  =+\infty
  \qquad\text{if }s(\rho)\not\leq s(\sigma).
\]
If $s(\rho)\leq s(\sigma)$, all the quantities may equivalently be
computed in the corner $P\mathcal B(\cH)P\cong\mathcal B(P\cH)$, where
$\sigma_P$ is faithful.  Matrices and states obtained in this corner
will always be regarded as matrices on $\cH$ by extension by zero on
$P^\perp\cH$.  In particular, under the support inclusion,
$Q_{\beta}^{(c)}(\rho\|\sigma)$ is finite and strictly positive for
every $\beta>0$, and the escort state
$\rho_\beta^{(c)}$ is well defined.




\begin{proposition}[Finite-dimensional fixed-ray representation]
\label{prop:fd-fixed-ray}
Let $\rho,\sigma$ be density matrices with
$s(\rho)\leq s(\sigma)$.  Let $\alpha,z>0$, $\alpha\ne1$, and put
$c=z/\alpha$.
Then
\begin{equation}
  D_{\alpha,z}(\rho\Vert\sigma)
  =\frac{\alpha}{\alpha-1}
   \int_1^\alpha
   \frac{D(\rho_{\beta}^{(c)}\Vert\sigma)}{\beta^2}
   \,\dd\beta .
  \label{eq:fd-fixed-ray-integral}
\end{equation}
The integral is oriented.  Thus, for $0<\alpha<1$,
\begin{equation}
  D_{\alpha,z}(\rho\Vert\sigma)
  =\frac{\alpha}{1-\alpha}
   \int_\alpha^1
   \frac{D(\rho_{\beta}^{(c)}\Vert\sigma)}{\beta^2}
   \,\dd\beta .
  \label{eq:fd-fixed-ray-lower}
\end{equation}
\end{proposition}

\begin{proof}
The calculation preceding the proposition proves the result when both
$\rho$ and $\sigma$ are faithful.  For general $\sigma$, put
$P=s(\sigma)$.  Since $s(\rho)\leq P$, we have $\rho=P\rho P$, and all
the moments, divergences, and escort states in the proposition agree
with the corresponding quantities computed in the corner
$\mathcal B(P\cH)$.  We may therefore replace $\cH$ by $P\cH$ and
$\sigma$ by its faithful restriction
$\sigma_P=\sigma\vert_{P\cH}$.

If $\rho$ is faithful on $P\cH$, the preceding calculation applies
directly.  Otherwise, define
\begin{equation}
  \rho_\varepsilon
  \defeq
  (1-\varepsilon)\rho
  +\varepsilon\frac{P}{\operatorname{rank}P},
  \qquad 0<\varepsilon<1.
  \label{eq:fd-faithful-regularization}
\end{equation}
Then $\rho_\varepsilon$ is a faithful density matrix on $P\cH$, and
hence
\begin{equation}
  D_{\alpha,z}(\rho_\varepsilon\Vert\sigma_P)
  =
  \frac{\alpha}{\alpha-1}
  \int_1^\alpha
  \frac{
    D((\rho_\varepsilon)_\beta^{(c)}\Vert\sigma_P)}
       {\beta^2}\,\dd\beta .
  \label{eq:fd-regularized-ray-integral}
\end{equation}

As $\varepsilon\downarrow0$, we have
$\rho_\varepsilon\to\rho$ in norm.  Continuity of the matrix power map
on the positive cone gives
\[
  \rho_\varepsilon^{1/c}\longrightarrow\rho^{1/c}
\]
in norm.  Consequently, for every fixed $\beta$ in the closed interval
with endpoints $1$ and $\alpha$,
\[
  X_\beta^{(c)}[\rho_\varepsilon]
  \longrightarrow
  X_\beta^{(c)}[\rho]
\]
and
\[
  \bigl(X_\beta^{(c)}[\rho_\varepsilon]\bigr)^{c\beta}
  \longrightarrow
  \bigl(X_\beta^{(c)}[\rho]\bigr)^{c\beta}
\]
in norm.  Taking traces therefore gives
\[
  Q_\beta^{(c)}(\rho_\varepsilon\|\sigma_P)
  \longrightarrow
  Q_\beta^{(c)}(\rho\|\sigma_P).
\]
The limiting moment is strictly positive: in the faithful corner the
powers of $\sigma_P$ are invertible, while $\rho^{1/c}\neq0$.  It
follows that
\[
  (\rho_\varepsilon)_\beta^{(c)}
  \longrightarrow
  \rho_\beta^{(c)}
\]
in norm for every such $\beta$.  In particular, at $\beta=\alpha$,
\[
  Q_{\alpha,z}(\rho_\varepsilon\|\sigma_P)
  \longrightarrow
  Q_{\alpha,z}(\rho\|\sigma_P)>0,
\]
and hence
\begin{equation}
  D_{\alpha,z}(\rho_\varepsilon\Vert\sigma_P)
  \longrightarrow
  D_{\alpha,z}(\rho\Vert\sigma_P).
  \label{eq:fd-regularized-divergence-limit}
\end{equation}

Since $\sigma_P$ is faithful and $P\cH$ is finite dimensional, the map
\[
  \eta\longmapsto D(\eta\Vert\sigma_P)
\]
is continuous on the compact state space of $\mathcal B(P\cH)$.
Therefore
\[
  D((\rho_\varepsilon)_\beta^{(c)}\Vert\sigma_P)
  \longrightarrow
  D(\rho_\beta^{(c)}\Vert\sigma_P)
\]
for every $\beta$.  Moreover, writing
$\lambda_{\min}(\sigma_P)>0$ for the smallest eigenvalue of
$\sigma_P$, every state $\eta$ on $P\cH$ satisfies
\begin{align}
  D(\eta\Vert\sigma_P)
  &=\Tr(\eta\log\eta)-\Tr(\eta\log\sigma_P) \notag\\
  &\leq-\log\lambda_{\min}(\sigma_P).
  \label{eq:fd-relative-entropy-uniform-bound}
\end{align}
Since $\beta^{-2}$ is bounded on the closed interval with endpoints
$1$ and $\alpha$, this gives an integrable bound independent of
$\varepsilon$.  Dominated convergence may therefore be applied to
\eqref{eq:fd-regularized-ray-integral}.  Combining it with
\eqref{eq:fd-regularized-divergence-limit} proves
\eqref{eq:fd-fixed-ray-integral}.  The lower-order form
\eqref{eq:fd-fixed-ray-lower} follows by reversing the orientation of
the integral.
\end{proof}

Denote the sandwiched R\'enyi divergence by $D_\alpha(\rho\|\sigma)$. At $c=1$, or equivalently $z=\alpha$, we write
\begin{align}
  x_\beta(\rho\|\sigma)
  &\defeq 
  \sigma^{\frac{1-\beta}{2\beta}}
  \rho
  \sigma^{\frac{1-\beta}{2\beta}},\\
  Q_\beta(\rho\|\sigma)
  &\defeq 
  \Tr x_\beta(\rho\|\sigma)^\beta,\\
  D_\beta(\rho\Vert\sigma)
  &\defeq 
  D_{\beta,\beta}(\rho\Vert\sigma),\\
  \rho_\beta^s
  &\defeq 
  \frac{x_\beta(\rho\|\sigma)^\beta}
       {Q_\beta(\rho\|\sigma)}.
\end{align}
Equation~\eqref{eq:fd-ray-refined-identity} becomes the refined
sandwiched identity of \cite{Bao:2019aol},
\begin{equation}
  \widetilde D_\beta(\rho\Vert\sigma)
  \defeq \beta^2\frac{\dd}{\dd\beta}
    \left[\frac{\beta-1}{\beta}
          D_\beta(\rho\Vert\sigma)\right]
  =D(\rho_\beta^s\Vert\sigma),
  \label{eq:fd-refined-srd}
\end{equation}
and Proposition~\ref{prop:fd-fixed-ray} gives
\begin{equation}
  D_\alpha(\rho\Vert\sigma)
  =\frac{\alpha}{\alpha-1}
   \int_1^\alpha
   \frac{D(\rho_\beta^s\Vert\sigma)}{\beta^2}\,\dd\beta .
  \label{eq:fd-srd-integral}
\end{equation}
This is an exact specialization at $z=\alpha$, not a limit in $z$.

\subsection{The lower-order support boundary}
\label{subsec:fd-support-boundary}

For $0<\a<1$, finiteness of {\az} divergence does not necessitate the support inclusion $s(\r)\leq s(\s)$.  The missing
support is then visible as an endpoint jump.  The finite-dimensional
formula in Proposition~\ref{prop:fd-boundary-correction} also clarifies why the general theorem in
Section~\ref{sec:results-and-consequences} assumes nested supports.

\begin{lemma}[Compression cost]\label{lem:compression}
Let $B\ge0$ be a positive semidefinite matrix on a finite-dimensional Hilbert space
$\mathcal H$, let $P\in\mathcal B(\mathcal H)$ be an orthogonal projection, and let $c>0$.
Then
\begin{equation}\label{eq:compression}
  0\ \le\ \operatorname{Tr}\big[(PBP)^{c}\big]\ \le\ \operatorname{Tr}\big[B^{c}\big],
\end{equation}
with equality on the left if and only if $s(B)\le 1-P$, and on the right if and only if
$s(B)\le P$.
\end{lemma}

\begin{proof}
Order eigenvalues of the compression being computed on $P\mathcal H$ decreasingly,
and extended by zeros. Denote the eigenvalues by $\lambda_i(PBP)$. For $v\in P\mathcal H$ one has $\langle v,PBPv\rangle=\langle v,Bv\rangle$,
so the Courant--Fischer-Weyl principle in the max-min form \cite[Corollary III.1.2]{Bhatia:1997}, applied over subspaces of $P\mathcal H$ rather than of $\mathcal H$, gives
\begin{equation}\label{eq:interlace}
  \lambda_{j}(PBP)\le\lambda_{j}(B)\qquad\text{for every }j.
\end{equation}
Since $t\mapsto t^{c}$ is increasing on $[0,\infty)$ and vanishes at $0$, summing
\eqref{eq:interlace} proves the right-hand inequality in \eqref{eq:compression}; the
left-hand one is trivial. (Only monotonicity is used, so the same argument gives
$\operatorname{Tr}f(PBP)\le\operatorname{Tr}f(B)$ for every increasing $f$ with $f(0)=0$.)

Suppose the right-hand equality holds. Each summand on the left is dominated by the
corresponding summand on the right, so all of them agree, and strict monotonicity of
$t\mapsto t^{c}$ upgrades this to equality in \eqref{eq:interlace} for every $j$. Summing
gives $\operatorname{Tr}[PBP]=\operatorname{Tr}[B]$, that is $\operatorname{Tr}[(1-P)B]=0$,
and $B\ge0$ forces $B^{1/2}(1-P)=0$, i.e.\ $s(B)\le P$. Conversely $s(B)\le P$ gives
$PBP=B$.

Finally, $\operatorname{Tr}[(PBP)^{c}]=0$ holds if and only if $(PBP)^{c}=0$, hence if and
only if $PBP=0$, hence if and only if $B^{1/2}P=0$, i.e.\ $s(B)\le 1-P$.
\end{proof}

\begin{proposition}[Finite-dimensional boundary correction]
\label{prop:fd-boundary-correction}
Let $0<\alpha<1$, $c>0$, and let $P=s(\sigma)$.  Define
\begin{equation}
  q_{1-,c}
  \defeq \Tr\!\left[(P\rho^{1/c}P)^c\right].
  \label{eq:fd-boundary-overlap}
\end{equation}
If $q_{1-,c}=0$, then
$D_{\alpha,c\alpha}(\rho\Vert\sigma)=+\infty$.  If $q_{1-,c}>0$, put
\begin{equation}
  \widehat\rho_c
  \defeq \frac{(P\rho^{1/c}P)^c}{q_{1-,c}}
  \quad\text{on }P\cH.
  \label{eq:fd-boundary-seed}
\end{equation}
Then, for every $0<\beta<1$, the fixed-ray escort states of $(\rho,\sigma)$ agree with those
of $(\widehat\rho_c,\sigma)$, and
\begin{equation}
  D_{\alpha,c\alpha}(\rho\Vert\sigma)
  =\frac{\alpha}{1-\alpha}
   \int_\alpha^1
   \frac{D((\widehat\rho_c)_\beta^{(c)}\Vert\sigma)}{\beta^2}
   \,\dd\beta
   -\frac{\alpha}{1-\alpha}\log q_{1-,c}.
  \label{eq:fd-alpha-c-alpha}
\end{equation}
For $c=1$, $q_{1-,1}=\Tr(P\rho)$. 

Moreover
\begin{equation}\label{eq:qbound}
  0\le q_{1-,c}\le1, 
\end{equation}
with
\begin{equation}\label{eq:endpoint-bounds}
    q_{1-,c}=1\iff s(\rho)\le s(\sigma), \qquad q_{1-,c}=0\iff s(\rho)\perp s(\sigma),
\end{equation}
so the boundary term in \eqref{eq:fd-alpha-c-alpha} is nonnegative and vanishes precisely under support inclusion.

Finally, the boundary correction arises as the limit
\begin{equation}\label{eq:endpointjump}
  \lim_{\beta\uparrow1}Q^{(c)}_{\beta}(\rho\|\sigma)=q_{1-,c}.
\end{equation}
\end{proposition}

\begin{proof}
The powers of $\sigma$ restrict the fixed-ray product to $P\cH$, so
\begin{equation}
  P\rho^{1/c}P=q_{1-,c}^{1/c}\widehat\rho_c^{1/c}.
\end{equation}
Consequently,
\begin{equation}
  Q_{\beta,c\beta}(\rho\|\sigma)
  =q_{1-,c}^{\beta}
   Q_{\beta,c\beta}(\widehat\rho_c\|\sigma),
  \qquad 0<\beta<1.
  \label{eq:fd-boundary-scaling}
\end{equation}
The scalar cancels in the normalized escort.  Taking logarithms at
$\beta=\alpha$ and applying Proposition~\ref{prop:fd-fixed-ray} to the
supported pair $(\widehat\rho_c,\sigma)$ proves
\eqref{eq:fd-alpha-c-alpha}.  
If $q_{1-,c}=0$, the lower-order moment
vanishes, and the convention $\log0=-\infty$ gives the asserted value.

Applying Lemma~\ref{lem:compression} with $B=\rho^{1/c}$, $P=s(\sigma)$ and using $\Tr(B^c)=\Tr(\rho)=1$ proves \eqref{eq:qbound} and \eqref{eq:endpoint-bounds}.

Finally, to prove the endpoint jump \eqref{eq:endpointjump} define
\[
    t_\beta\defeq\frac{1-\beta}{2c \beta}.
\]
As $t_\beta\longrightarrow0$, $\sigma_P^{t_\beta}\longrightarrow P $ in norm. Since, $c\beta \longrightarrow c>0$, joint norm continuity of $(A,t)\longmapsto A^t$ for positive finite dimensional matrices $A$ and exponents bounded away from zero gives
\[
   \left( X_{\beta}^{(c)}\right)^{c\beta}\longrightarrow (P \rho^{1/c} P)^c.
\]
Taking traces proves
\[
    \lim_{\beta\uparrow1} Q_\beta^{(c)}(\rho\|\sigma)= \Tr((P \rho^{1/c} P)^c)=q_{1-,c}. \qedhere
\]
\end{proof}

\begin{example}[Why the boundary term cannot be omitted]
\label{ex:two-point-boundary}
Let $\rho=(r,1-r)$ and $\sigma=(1,0)$ on $\mathbb C^2$, where
$0<r<1$.  On the sandwiched ray $c=1$, every interior escort with
$0<\beta<1$ equals $\sigma$.  The escort integral therefore vanishes, but
\begin{equation}
  D_\alpha(\rho\Vert\sigma)
  =-\frac{\alpha}{1-\alpha}\log r,
  \qquad 0<\alpha<1.
\end{equation}
This is precisely the support-boundary term with
$r=\rho(s(\sigma))$.
\end{example}

\section{Operator-algebraic preliminaries}
\label{sec:preliminaries}

We prove the integral formulas for general von Neumann algebras in the next section. Here, we collect only the material needed to state the general results.  Standard
references for Haagerup and Kosaki spaces include
\cite{Terp:1981lp,KOSAKI198429}, and the particular $\alpha$--$z$ construction used
here follows \cite{Kato:2023aro,Kato:2023hlj,Hiai:2024qve}. We also refer to \cite{Berta:2016vnw,Jencova:2016tqz,Jencova:2017txf} for sandwiched R\'enyi divergences.

\subsection{Haagerup \texorpdfstring{$L^p$}{Lp} spaces, densities, and support corners}

Let $M$ be a von Neumann algebra and choose a faithful normal semifinite
weight.  Its continuous core $\cN$ carries a canonical semifinite trace
$\tau$ and a dual action $(\theta_s)_{s\in\mathbb R}$.  For $0<p<\infty$,
the Haagerup space $L^p(M)$ consists of the $\tau$-measurable operators
affiliated with $\cN$ that satisfy
\begin{equation}
  \theta_s(x)=\e^{-s/p}x,
  \qquad s\in\mathbb R.
  \label{eq:haagerup-homogeneity}
\end{equation}
The construction is independent, up to canonical isometry, of the initial
weight.  For $p\geq1$ it is a Banach space; for $0<p<1$ it is a complete
quasi-Banach space. The standard $L^p(M)$ (quasi-)norm is denoted by
$\|\cdot\|_p$. Multiplication is understood as closed multiplication
of measurable operators and obeys the generalized H\"older inequality.

There is a canonical isometric identification
$M_*\cong L^1(M)$.  In particular, every
$\varphi\in M_*^+$ has a unique density
$h_\varphi\in L^1(M)_+$.
Moreover, there is a canonical functional $\tr$ on $L^1(M)$ such that
\begin{equation}
  \tr(h_\varphi)=\varphi(\one),
  \label{eq:density-normalization}
\end{equation}
and
\begin{equation}
    \tr(|h_{\varphi}|)=\tr(h_{|\varphi|})=|\varphi|(\one)=\|\varphi\|.
\end{equation}
The support of $h_\varphi$ corresponds to the support projection
$s(\varphi)\in M$.  For a projection $e\in M$, there is a canonical
identification
\begin{equation}
  L^p(eMe)\cong eL^p(M)e.
  \label{eq:corner-identification}
\end{equation}
In particular, after replacing $M$ by
$eMe$, $e=s(\omega)$, a normal state $\omega$ becomes faithful.  Since a
faithful normal state exists on this corner, no global
$\sigma$-finiteness hypothesis on $M$ is needed in our main results.

We write $D(\psi\Vert\omega)$ for Araki relative entropy of normal states,
with value in $[0,+\infty]$.  We extend the first argument to a nonzero
normal positive functional $\nu=q\varphi$, where $q=\nu(\one)>0$ and
$\varphi$ is a state, by the standard scaling law
\begin{equation}
  D(q\varphi\Vert\omega)
  \defeq qD(\varphi\Vert\omega)+q\log q.
  \label{eq:araki-relative-entropy-scaling}
\end{equation}
Its normalized relative entropy is therefore
\begin{equation}
  D_1(\nu\Vert\omega)
  \defeq \frac{D(\nu\Vert\omega)}{\nu(\one)}
  =D(\varphi\Vert\omega)+\log q.
  \label{eq:relative-entropy-scaling}
\end{equation}
Here and below the second argument is a state.

\subsection{Algebraic \texorpdfstring{$\alpha$--$z$}{alpha-z} and
sandwiched R\'enyi divergences}

Let $\psi$ and $\omega$ be normal states, and put $e=s(\omega)$.  For
$0<\alpha<1$ and $z>0$, define the closed positive product
\begin{equation}
  x_{\alpha,z}(\psi\|\omega)
  \defeq \overline{
    h_\omega^{\frac{1-\alpha}{2z}}
    h_\psi^{\frac\alpha z}
    h_\omega^{\frac{1-\alpha}{2z}}
  }\in eL^z(M)_+e
  \label{eq:lower-alpha-z-factor}
\end{equation}
and
\begin{equation}
  Q_{\alpha,z}(\psi\|\omega)
  \defeq \tr\bigl(x_{\alpha,z}(\psi\|\omega)^z\bigr).
  \label{eq:lower-alpha-z-moment}
\end{equation}
In more detail, set
\begin{equation}
    A = h_\omega^{\frac{1-\alpha}{2z}},
    \qquad
    C = h_\psi^{\frac\alpha{2z}},
\end{equation}
and define
\begin{equation}
    T_{\alpha,z}
    \defeq \overline{CA},
    \qquad
    x_{\alpha,z}
    =T_{\alpha,z}^*T_{\alpha,z}.
\end{equation}
Equivalently, in the associative $*$-algebra of $\tau$-measurable operators, $x_{\alpha,z}$ is the positive self-adjoint product
\begin{equation}
    x_{\alpha,z}
    =T_{\alpha,z}^*T_{\alpha,z}
    =\overline{AC^2A}
    =\overline{
    h_\omega^{\frac{1-\alpha}{2z}}
    h_\psi^{\frac\alpha z}
    h_\omega^{\frac{1-\alpha}{2z}}
  }.
\end{equation}
For $\alpha>1$, the moment $Q_{\alpha,z}(\psi\|\omega)$ is finite when there is a positive factor \cite{Kato:2023aro, Kato:2023hlj}
\begin{equation}
  x_{\alpha,z}(\psi\|\omega)\in eL^z(M)_+e
  \label{eq:upper-alpha-z-corner}
\end{equation}
satisfying
\begin{equation}
  h_\psi^{\alpha/z}
  =h_\omega^{\frac{\alpha-1}{2z}}
   x_{\alpha,z}(\psi\|\omega)
   h_\omega^{\frac{\alpha-1}{2z}}.
  \label{eq:upper-alpha-z-factor}
\end{equation}
The factor in the support corner is unique.  In this case
\begin{equation}
  Q_{\alpha,z}(\psi\|\omega)
  \defeq \tr\bigl(x_{\alpha,z}(\psi\|\omega)^z\bigr);
  \label{eq:upper-alpha-z-moment}
\end{equation}
otherwise it is $+\infty$.  The {\az} divergence is
\begin{equation}
  D_{\alpha,z}(\psi\Vert\omega)
  \defeq \frac{1}{\alpha-1}
  \log Q_{\alpha,z}(\psi\|\omega).
  \label{eq:algebraic-alpha-z-divergence}
\end{equation}

The same factor and moment definitions apply to a nonzero normal positive
first argument $\nu$ after replacing $h_\psi$ by $h_\nu$.  In particular,
for $q>0$ and a state $\varphi$,
\begin{equation}
  Q_{\alpha,z}(q\varphi\|\omega)
  =q^\alpha Q_{\alpha,z}(\varphi\|\omega).
  \label{eq:alpha-z-moment-homogeneity}
\end{equation}
For such an argument we employ the normalized convention
\begin{equation}
  D_{\alpha,z}(\nu\Vert\omega)
  \defeq \frac{1}{\alpha-1}
  \log\frac{Q_{\alpha,z}(\nu\|\omega)}{\nu(\one)}.
  \label{eq:normalized-alpha-z-positive}
\end{equation}
For every fixed $z>0$, its lower-order limit is
\begin{equation}
  \lim_{\alpha\uparrow1}D_{\alpha,z}(\nu\Vert\omega)
  =D_1(\nu\Vert\omega).
  \label{eq:lower-alpha-z-limit}
\end{equation}
The upper-order limit $\alpha\downarrow1$ is known under the additional
conditions $z>1/2$ and
$D_{\alpha_0,z}(\nu\Vert\omega)<\infty$ for some
$\alpha_0\in(1,2z]$ \cite[Theorems~6.7 and~6.9]{Hiai:2024qve}.

The sandwiched quantities are the specialization
\begin{equation}
  Q_\alpha(\psi\|\omega)
  \defeq Q_{\alpha,\alpha}(\psi\|\omega),
  \qquad
  D_\alpha(\psi\Vert\omega)
  \defeq D_{\alpha,\alpha}(\psi\Vert\omega).
  \label{eq:srd-specialization}
\end{equation}
For $0<\alpha<1$ this means
\begin{equation}
  x_\alpha^s(\psi\|\omega)
  \defeq \overline{
  h_\omega^{\frac{1-\alpha}{2\alpha}}
  h_\psi
  h_\omega^{\frac{1-\alpha}{2\alpha}}},
  \qquad
  Q_\alpha(\psi\|\omega)
  =\tr\bigl(x_\alpha^s(\psi\|\omega)^\alpha\bigr).
  \label{eq:lower-srd-factor}
\end{equation}
For $\alpha>1$ with finite moment, let
$x_\alpha^s(\psi\|\omega)\in eL^\alpha(M)_+e$ be the upper factor
satisfying
\begin{equation}
  h_\psi
  =h_\omega^{\frac{\alpha-1}{2\alpha}}
   x_\alpha^s(\psi\|\omega)
   h_\omega^{\frac{\alpha-1}{2\alpha}},
  \qquad
  Q_\alpha(\psi\|\omega)
  =\tr\bigl(x_\alpha^s(\psi\|\omega)^\alpha\bigr).
  \label{eq:upper-srd-factor}
\end{equation}
At the endpoint we set
\begin{equation}
  x_1^s(\psi\|\omega)\defeq h_\psi.
  \label{eq:srd-factor-at-one}
\end{equation}

The full data-processing region for $D_{\alpha,z}$ is \cite{Zhang:2018roy, Hiai:2024qve}
\begin{equation}
  \begin{cases}
  0<\alpha<1,
  &z\geq\max\{\alpha,1-\alpha\},\\[0.4ex]
  \alpha>1,
  &\max\{\alpha/2,\alpha-1\}\leq z\leq\alpha.
  \end{cases}
  \label{eq:dpi-region}
\end{equation}
This restriction will be imposed only for
consequences that use data processing, not for the integral theorem.

\subsection{Fixed-ray factors and escort trajectories}
\label{subsec:algebraic-ray-escorts}

Fix $c>0$ and keep $e=s(\omega)$.  For $0<\beta<1$, define
\begin{equation}
  x_\beta^{(c)}(\psi\|\omega)
  \defeq \overline{
  h_\omega^{\frac{1-\beta}{2c\beta}}
  h_\psi^{1/c}
  h_\omega^{\frac{1-\beta}{2c\beta}}}
  \in eL^{c\beta}(M)_+e.
  \label{eq:lower-fixed-ray-factor}
\end{equation}
At $\beta=1$ set
$x_1^{(c)}=h_\psi^{1/c}$ and
$Q_1^{(c)}(\psi\|\omega)=Q_{1,c}(\psi\|\omega)\defeq 1$.  If
$s(\psi)\nleq e$, this is a formal endpoint outside the corner; it belongs
to $eL^c(M)e$ precisely for a supported pair.  This is the endpoint jump
measured by the boundary terms below.  
For $\beta>1$, whenever it exists, let
$x_\beta^{(c)}\in eL^{c\beta}(M)_+e$ be the unique positive factor such that
\begin{equation}
  h_\psi^{1/c}
  =h_\omega^{\frac{\beta-1}{2c\beta}}
   x_\beta^{(c)}
   h_\omega^{\frac{\beta-1}{2c\beta}}.
  \label{eq:upper-fixed-ray-factor}
\end{equation}
In both cases put
\begin{equation}
  Q_\beta^{(c)}(\psi\|\omega)
  \defeq \tr\bigl((x_\beta^{(c)})^{c\beta}\bigr)
  =Q_{\beta,c\beta}(\psi\|\omega),
  \label{eq:algebraic-ray-moment}
\end{equation}
with $ Q_\beta^{(c)}(\psi\|\omega)=+\infty$ when the factor $x_\beta^{(c)}$ does not exist.
Whenever this number is finite and nonzero, the fixed-ray escort state $\psi_\beta^{(c)}$ is defined in terms of its associated $L^1$ density
\begin{equation}
  h_{\psi_\beta^{(c)}}
  \defeq \frac{(x_\beta^{(c)})^{c\beta}}
          {Q_\beta^{(c)}(\psi\|\omega)} \ .
  \label{eq:algebraic-ray-escort}
\end{equation}
At $c=1$ these are precisely the sandwiched factors and escorts.  We set
$\psi_\beta^s\defeq \psi_\beta^{(1)}$ and have
\begin{equation}
  x_\beta^{(1)}=x_\beta^s,
  \qquad
  Q_\beta^{(1)}=Q_\beta,
  \qquad \psi_\beta^{(1)}=\psi_\beta^s.
  \label{eq:sandwiched-ray-specialization}
\end{equation}

\section{Main results and consequences}
\label{sec:results-and-consequences}

We now state the integral representations for arbitrary von Neumann
algebras and collect their consequences.  Their proofs are postponed to
Sections~\ref{sec:proof-lower-srd} and \ref{sec:proof-alpha-z}.

For $\alpha>0$, $\alpha\ne1$, define
\begin{equation}
  I_\alpha
  \defeq \begin{cases}
     [1,\alpha],&\alpha>1,\\
     [\alpha,1],&0<\alpha<1,
     \end{cases}
  \label{eq:order-interval}
\end{equation}
and the probability measure
\begin{equation}
  \dd\mu_\alpha(\beta)
  \defeq \begin{cases}
  \displaystyle
  \frac{\alpha}{\alpha-1}\frac{\dd\beta}{\beta^2},
  &\alpha>1,\\[1.2ex]
  \displaystyle
  \frac{\alpha}{1-\alpha}\frac{\dd\beta}{\beta^2},
  &0<\alpha<1.
  \end{cases}
  \label{eq:order-probability-measure}
\end{equation}

For clarity, Table~\ref{tab:scope} separates the analytic theorem from
the narrower parameter ranges imposed by data processing or by an imported
operational result.

\begin{table}[ht]
\centering
\small
\begin{tabular}{L{0.23\textwidth}L{0.29\textwidth}L{0.31\textwidth}}
\hline
Result & Parameters & Algebra;  additional hypotheses \\
\hline
Fixed-ray integral, lower
& $0<\alpha<1$, $z>0$
& Arbitrary; $s(\psi)\leq s(\omega)$ \\
&&\\
Fixed-ray integral, upper
& $\alpha>1$, $z>0$
& Arbitrary; $Q_{\alpha,z}(\psi\|\omega)<\infty$ \\
&&\\
Boundary-corrected integral
& $0<\alpha<1$, $z>0$
& Arbitrary; no support inclusion \\
&&\\
One-shot testing
& upper part of \eqref{eq:dpi-region}
& Arbitrary \\
&&\\
Strong converse
& sandwiched orders $\alpha>1$
& Hypotheses of the cited operational theorems \\
&&\\
Work reliability
& sandwiched orders $0<\alpha<1$
& finite dimensional; thermal operations and the battery model of
  \cite{Watanabe:2026pzw} \\
&&\\
Pair conversion
& $0<\alpha<1$, $z>\max\{\alpha,1-\alpha\}$
& finite dimensional; restricted class $\mathcal F$ and source conditions \\
\hline
\end{tabular}
\caption{Parameter ranges and hypotheses. ``Arbitrary'' here means an arbitrary
von Neumann algebra.}
\label{tab:scope}
\end{table}

\subsection{Integral theorems and the support boundary}

\begin{theorem}[Fixed-ray $\alpha$--$z$ integral representation]
\label{thm:fixed-ray-integral}
Let $M$ be an arbitrary von Neumann algebra and let $\psi,\omega$ be
normal states.  Fix $\alpha,z>0$, $\alpha\ne1$, and put $c=z/\alpha$.
Assume, branchwise, that
\begin{equation}
  \begin{cases}
  s(\psi)\leq s(\omega),&0<\alpha<1,\\
  Q_{\alpha,z}(\psi\|\omega)<\infty,&\alpha>1.
  \end{cases}
  \label{eq:main-branchwise-hypotheses}
\end{equation}
In the upper branch, finiteness already implies support inclusion.  Then,
for every $\beta\in I_\alpha$, the fixed-ray factor $x_\beta^{(c)}$ exists
and its moment is finite and nonzero, so that
\eqref{eq:algebraic-ray-escort} defines a normal state.  Moreover,
\begin{equation}
  D_{\alpha,z}(\psi\Vert\omega)
  =\frac{\alpha}{\alpha-1}
   \int_1^\alpha
   \frac{D(\psi_\beta^{(c)}\Vert\omega)}{\beta^2}
   \,\dd\beta
  =\int_{I_\alpha}
   D(\psi_\beta^{(c)}\Vert\omega)\,
   \dd\mu_\alpha(\beta).
  \label{eq:main-fixed-ray-integral}
\end{equation}
The first integral is oriented.  Both integrals are Lebesgue integrals
determined by the open interval between $1$ and $\alpha$; consequently, an
infinite relative-entropy value at an endpoint does not affect them.
\end{theorem}

The theorem is stated on the original algebra, but its construction may be
performed after passing to the support corner
$eMe$, $e=s(\omega)$.  There $\omega$ is faithful and the corner is
$\sigma$-finite.  Thus neither faithfulness on the original algebra nor
global $\sigma$-finiteness is an additional hypothesis.

There is also a sharp lower-order formula without support inclusion.  
\begin{proposition}[General lower-order fixed-ray boundary]
\label{prop:general-fixed-ray-boundary}
Let $M$ be a von Neumann algebra, let $\psi,\omega$ be normal states, let
$0<\alpha<1$, and let $c>0$.  Put
$e=s(\omega)$, $\omega_e\defeq \omega\vert_{eMe}$, and, under the canonical
identification $L^c(eMe)\cong eL^c(M)e$, define
\begin{equation}
  y_c\defeq \overline{e h_\psi^{1/c}e}\in L^c(eMe)_+,
  \qquad
  q_c\defeq \tr_{eMe}(y_c^c).
  \label{eq:general-boundary-anchor}
\end{equation}
If $q_c=0$, then
$D_{\alpha,c\alpha}(\psi\Vert\omega)=+\infty$.  If $q_c>0$, define a
normal state $\widehat\psi_c$ on $eMe$ by
\begin{equation}
  h_{\widehat\psi_c}\defeq q_c^{-1}y_c^c.
  \label{eq:general-boundary-seed}
\end{equation}
Then, for every $0<\beta<1$, the fixed-ray escort of $(\psi,\omega)$ is
the canonical extension to $M$ of the fixed-ray escort of
$(\widehat\psi_c,\omega_e)$, and
\begin{equation}
  D_{\alpha,c\alpha}(\psi\Vert\omega)
  =\frac{\alpha}{1-\alpha}
   \int_\alpha^1
   \frac{D((\widehat\psi_c)_\beta^{(c)}\Vert\omega_e)}{\beta^2}
   \,\dd\beta
   -\frac{\alpha}{1-\alpha}\log q_c.
  \label{eq:general-fixed-ray-boundary}
\end{equation}
Moreover,
\begin{equation}\label{eq:boundary-cost-bound}
    0\leq q_c \leq 1.
\end{equation}
\end{proposition}

\begin{corollary}[Sandwiched integral representation]
\label{cor:srd-integral}
Under the hypotheses of Theorem~\ref{thm:fixed-ray-integral}, set
$z=\alpha$, so that $c=1$.  Then
\begin{equation}
  D_\alpha(\psi\Vert\omega)
  =\frac{\alpha}{\alpha-1}
   \int_1^\alpha
   \frac{D(\psi_\beta^s\Vert\omega)}{\beta^2}\,\dd\beta .
  \label{eq:main-srd-integral}
\end{equation}
Analytically, this holds for every $0<\alpha<1$ under
$s(\psi)\leq s(\omega)$, and for every
$\alpha>1$ for which $Q_\alpha(\psi\|\omega)<\infty$.
\end{corollary}

For lower orders, the nested-support assumption in
Corollary~\ref{cor:srd-integral} is essential.  The general statement for
normal states is the following.

\begin{corollary}[Support-boundary correction for lower-order SRD]
\label{cor:srd-support-boundary}
Let $M$ be a von Neumann algebra, let $\psi,\omega$ be normal states, and
let $0<\alpha<1$.  Put
\begin{equation}
  e\defeq s(\omega),\qquad q\defeq \psi(e).
  \label{eq:srd-overlap}
\end{equation}
If $q=0$, then
$Q_\alpha(\psi\|\omega)=0$ and
$D_\alpha(\psi\Vert\omega)=+\infty$.  If $q>0$, let
$\omega_e=\omega\vert_{eMe}$ and define the normal state
$\widehat\psi$ on $eMe$ by
\begin{equation}
  \widehat\psi(x)\defeq q^{-1}\psi(x),
  \qquad x\in eMe.
  \label{eq:normalized-compression}
\end{equation}
Then
\begin{equation}
  D_\alpha(\psi\Vert\omega)
  =\frac{\alpha}{1-\alpha}
   \int_\alpha^1
   \frac{D(\widehat\psi_\beta^s\Vert\omega_e)}{\beta^2}
   \,\dd\beta
   -\frac{\alpha}{1-\alpha}\log q.
  \label{eq:srd-support-boundary}
\end{equation}
For every $0<\beta<1$, the escort constructed from $(\psi,\omega)$ is the
canonical extension of $\widehat\psi_\beta^s$, namely
$x\mapsto\widehat\psi_\beta^s(exe)$.  The boundary term vanishes exactly
when $s(\psi)\leq s(\omega)$.
\end{corollary}

% At $c=1$, $q_1=\tr(eh_\psi e)=\psi(e)$, so
% Corollary~\ref{cor:srd-support-boundary} is precisely the specialization
% of Proposition~\ref{prop:general-fixed-ray-boundary}.  In finite dimensions,
% $q_c=\Tr[(P\rho^{1/c}P)^c]$, recovering
% Proposition~\ref{prop:fd-boundary-correction}.

The proofs of the results stated above are collected in Section~\ref{sec:proof-alpha-z}. Section~\ref{subsec:compatible-scale} constructs a compatible weighted Haagerup scale and establishes the anchored entropy identity that provides the analytic input. Section~\ref{subsec:fixed-ray-proof} identifies the fixed-ray factors and escort states with the corresponding anchored path and proves Theorem~\ref{thm:fixed-ray-integral}; Corollary~\ref{cor:srd-integral} then follows by setting \(c=1\). Finally, Section~\ref{subsec:lower-order-proof} proves the bound on \(q_c\), the escort identification, and the boundary formula of Proposition~\ref{prop:general-fixed-ray-boundary}, including the case \(q_c=0\), with Corollary~\ref{cor:srd-support-boundary} following as its specialization to \(c=1\).

\subsection{Refinement, monotonicity, and convex order geometry}

For a support-nested pair of states $\psi,\omega$ with $s(\psi)\leq s(\omega)$, and fixed $c>0$, the map
$\beta\mapsto\psi_\beta^{(c)}$ is the \emph{fixed-ray escort trajectory}.
Use the endpoint convention $Q_{1,c}=Q_1^{(c)}=1$ from
Section~\ref{subsec:algebraic-ray-escorts}, and define
\begin{equation}
  \begin{gathered}
  \mathcal J_c
  \defeq\{\beta>0:0<Q_{\beta,c\beta}(\psi\|\omega)<\infty\},\\
  R_c(\beta)\defeq D(\psi_\beta^{(c)}\Vert\omega),
  \qquad
  F_c(\beta)\defeq\frac1\beta
       \log Q_{\beta,c\beta}(\psi\|\omega),
  \qquad \beta\in\mathcal J_c.
  \end{gathered}
  \label{eq:escort-profile}
\end{equation}
For $c=1$, $R_1(\beta)$ is the object referred to as refined SRD in \cite{Bao:2019aol}. We will sometimes refer to $R_c(\beta)$ as \textit{fixed-ray refined divergence}. 

Supportedness $s(\psi)\leq s(\omega)$ gives $(0,1]\subseteq\mathcal J_c$ using Theorem~\ref{thm:fixed-ray-integral}.  If
$\beta_2>1$ belongs to $\mathcal J_c$, Theorem~\ref{thm:fixed-ray-integral}
applied with endpoint $\beta_2$ shows that
$[1,\beta_2]\subseteq\mathcal J_c$.  Hence $\mathcal J_c$ is an interval.
The proof of Theorem~\ref{thm:fixed-ray-integral} in
Section~\ref{subsec:fixed-ray-proof} also yields the following raywise
monotonicity and derivative identity. Let $\frac{\dd^-}{\dd\beta}$ denote the left derivative with respect to $\beta$. 

\begin{corollary}[Raywise refinement and monotonicity]
\label{cor:raywise-monotonicity}
The function $R_c$ is nondecreasing on $\mathcal J_c$.  At every interior
point of $\mathcal J_c$,
\begin{equation}
  \beta^2\frac{\dd^-}{\dd\beta}F_c(\beta)
  =R_c(\beta).
  \label{eq:raywise-refined-identity}
\end{equation}
At $\beta=1$ the same identity holds with the left derivative at the
endpoint, even when $\mathcal J_c=(0,1]$; if
$D(\psi\Vert\omega)=+\infty$, this derivative is understood in the
extended-real sense.  For $\beta\neq1$,
\[
  F_c(\beta)=\frac{\beta-1}{\beta}
  D_{\beta,c\beta}(\psi\Vert\omega),
  \qquad F_c(1)=0.
\]
In finite dimensions these functions are smooth on the interior of
$\mathcal J_c$, and the left derivative there may be replaced by the
ordinary derivative.
\end{corollary}

Thus the finite-dimensional refined divergence of \cite{Bao:2019aol} extends intrinsically as a
one-sided slope along the escort trajectory.  Since $\mu_\alpha$ is a probability measure, for every $\alpha\in \mathcal J_c\backslash\{1\} $, monotonicity
also gives the endpoint bounds
\begin{align}
  D(\psi\Vert\omega)
  &\leq D_{\alpha,c\alpha}(\psi\Vert\omega)
  \leq D(\psi_\alpha^{(c)}\Vert\omega),
  &&\alpha>1,
  \label{eq:upper-endpoint-bounds}\\
  D(\psi_\alpha^{(c)}\Vert\omega)
  &\leq D_{\alpha,c\alpha}(\psi\Vert\omega)
  \leq D(\psi\Vert\omega),
  &&0<\alpha<1,
  \label{eq:lower-endpoint-bounds}
\end{align}
in the extended sense.

It is useful to reparametrize the order by
\begin{equation}
  \kappa\defeq 1-\frac1\alpha,
  \qquad
  \alpha=\frac1{1-\kappa},
  \qquad -\infty<\kappa<1.
  \label{eq:kappa-parameter}
\end{equation}
Define the effective finite-moment domain
\begin{equation}
  \mathcal K_c
  \defeq \left\{\kappa<1:
  0<Q_{\frac1{1-\kappa},\frac{c}{1-\kappa}}
       (\psi\|\omega)<\infty\right\}
  =\left\{\kappa<1:\frac1{1-\kappa}\in\mathcal J_c\right\}.
  \label{eq:ray-effective-domain}
\end{equation}
Since $\kappa\mapsto1/(1-\kappa)$ is increasing and $\mathcal J_c$ is an
interval, $\mathcal K_c$ is an interval.  For a support-nested pair of states,
$(-\infty,0]\subseteq\mathcal K_c$.
Define the fixed-ray potential on
$\mathcal K_c$ by
\begin{equation}
  \Phi_c(\kappa)
  \defeq \kappa
  D_{\frac1{1-\kappa},\frac{c}{1-\kappa}}
       (\psi\Vert\omega),
  \qquad \Phi_c(0)\defeq 0,
  \label{eq:ray-potential}
\end{equation}
and extend it by $+\infty$ on $(-\infty,1)\setminus\mathcal K_c$ if an
extended-real-valued convex function is desired.
Changing variables $u=1-1/\beta$ in
\eqref{eq:main-fixed-ray-integral} gives
\begin{equation}
  \Phi_c(\kappa)
  =\int_0^\kappa
   R_c\!\left(\frac1{1-u}\right)\,\dd u,
  \label{eq:ray-potential-integral}
\end{equation}
for every $\kappa\in\mathcal K_c$.  As an integral of the nondecreasing function $u\mapsto R_c\left(\frac1{1-u} \right)$, $\Phi_c$ is convex on
its effective domain.  For the formal convex transform
\begin{equation}
  E_c(r)
  \defeq 
  \sup_{\kappa\in\mathcal K_c\cap[0,1)} \{\kappa r-\Phi_c(\kappa)\}
  =
  \max\!\left\{
    0,\,
    \sup_{\substack{\alpha>1\\
    0<Q_{\alpha,c\alpha}(\psi\Vert\omega)<\infty}}
    \frac{\alpha-1}{\alpha}
    \bigl(r-D_{\alpha,c\alpha}(\psi\Vert\omega)\bigr)
  \right\}.
  \label{eq:formal-convex-transform}
\end{equation}
The inclusion of $\kappa=0$ ensures that $E_c(r)\geq0$, even when no
upper-order moment is finite.  If an optimizer $\kappa_r>0$ lies in the
interior of the effective domain, and $\alpha_r=1/(1-\kappa_r)$, define
$R_c(\alpha_r\pm)$ as the corresponding one-sided limits within
$\operatorname{int}\mathcal J_c$.  The subgradient optimality condition is
\begin{equation}
  R_c(\alpha_r-)
  \leq r\leq R_c(\alpha_r+).
  \label{eq:nonsmooth-level-crossing}
\end{equation}
At a continuity point this is the level-crossing equation
\begin{equation}
  D(\psi_{\alpha_r}^{(c)}\Vert\omega)=r.
  \label{eq:smooth-level-crossing}
\end{equation}
For the sandwiched ray \(c=1\), the transform \(E_1\) is identified,
under the hypotheses of the relevant operational theorems, with the exact
strong-converse exponent of asymmetric binary hypothesis testing; see
Section~\ref{subsec:binary-testing}.
For general \(c\), \(E_c\) remains an intrinsic convex transform unless a
separate operational identification is available.



\subsection{Stability and the escort average}

Assume the branchwise hypotheses of
Theorem~\ref{thm:fixed-ray-integral}: for \(0<\alpha<1\), assume
\(s(\psi)\leq s(\omega)\), whereas for \(\alpha>1\), assume
\(Q_{\alpha,z}(\psi\|\omega)<\infty\).  Put \(c=z/\alpha\) and write
\(\|\cdot\|_1=\|\cdot\|_{M_*}\) for the predual norm.

Pinsker's inequality for Araki relative entropy \cite[Theorem~3.1]{Hiai1981SufficiencyKC} gives, pointwise,
\begin{equation}
  R_c(\beta)
  =
  D(\psi_\beta^{(c)}\Vert\omega)
  \geq
  \frac12
  \|\psi_\beta^{(c)}-\omega\|_1^2,
  \qquad \beta\in I_\alpha .
  \label{eq:escort-pinsker-pointwise}
\end{equation}
We emphasize that this statement requires neither finite dimensionality nor any
measurability assumption on the state-valued escort path. The result in \cite{Hiai1981SufficiencyKC} assumes a cyclic and separating vector. This directly covers our setting if $\omega$ is faithful normal. For arbitrary $M$, the inequality follows by restricting to the $\sigma$-finite corner $eMe$ with $e=s(\psi_\beta^{(c)}+\omega)$, where the relative entropy and predual distance are unchanged.

To formulate the averaged inequality without imposing measurability, define the upper Lebesgue integral of a nonnegative function
\(f\) by
\begin{equation}
  \int_{I_\alpha}^{*} f\,\dd\mu_\alpha
  :=
  \inf\left\{
    \int_{I_\alpha}g\,\dd\mu_\alpha:
    g\text{ is Borel measurable and }f\leq g
  \right\}.
  \label{eq:upper-lebesgue-integral}
\end{equation}
The proof of Corollary~\ref{cor:raywise-monotonicity} in Section \ref{sec:proof-alpha-z} shows that $R_c$ is
Borel measurable.
Using this together with
\(\|\psi_\beta^{(c)}-\omega\|_1^2\leq2R_c(\beta)\),
Theorem~\ref{thm:fixed-ray-integral} implies
\begin{equation}
  D_{\alpha,z}(\psi\Vert\omega)
  \geq
  \frac12
  \int_{I_\alpha}^{*}
    \|\psi_\beta^{(c)}-\omega\|_1^2
    \,\dd\mu_\alpha(\beta).
  \label{eq:escort-pinsker-upper-integral}
\end{equation}
Thus \eqref{eq:escort-pinsker-upper-integral} holds on an arbitrary von
Neumann algebra, without any data-processing, faithfulness,
\(\sigma\)-finiteness, or finite-dimensionality assumption beyond the
hypotheses of the fixed-ray integral theorem.

If the scalar function
\[
  \beta\longmapsto
  \|\psi_\beta^{(c)}-\omega\|_1^2
\]
is \(\mu_\alpha\)-measurable, then the upper integral is the ordinary
Lebesgue integral and
\begin{equation}
  D_{\alpha,z}(\psi\Vert\omega)
  \geq
  \frac12
  \int_{I_\alpha}
    \|\psi_\beta^{(c)}-\omega\|_1^2
    \,\dd\mu_\alpha(\beta).
  \label{eq:escort-pinsker-ordinary}
\end{equation}
Finite dimensionality is one sufficient condition for this scalar
measurability, since the escort path is then trace-norm continuous, but it
is not necessary.

For $\varepsilon>0$, let
\begin{equation}
  f(\beta)= \|\psi_\beta^{(c)}-\omega\|_1^2, \quad A_\varepsilon  \defeq \{f\geq \varepsilon^2\},
\end{equation}
and let $\mu_\alpha^*$ be the outer measure induced by $\mu_\alpha$.\footnote{If $\mu$ is a measure on $(X,\Sigma)$, its induced outer measure is the function $\mu^*:2^X\to[0,\infty]$ defined by $\mu^*(A)=\inf\{\mu(B):B\in\Sigma,A\subseteq B\}$, where $2^X$ denotes the collection of all subsets of $X$.}
Then for every Borel $g\geq f$,
\begin{equation}
    \mu_\alpha^*(A_\varepsilon)\leq \mu_\alpha\{g\geq \varepsilon^2\}\leq \frac{1}{\varepsilon^2} \int g \, d\mu_\alpha.
\end{equation}
Taking the infimum over such $g$ and using \eqref{eq:escort-pinsker-upper-integral} gives, for every \(\varepsilon>0\),
\begin{equation}
  \mu_\alpha^*\!\left(
    \left\{\beta\in I_\alpha:
      \|\psi_\beta^{(c)}-\omega\|_1\geq\varepsilon
    \right\}
  \right)
  \leq
  \min\left\{
    1,\frac{2D_{\alpha,z}(\psi\Vert\omega)}{\varepsilon^2}
  \right\},
  \label{eq:escort-concentration}
\end{equation}
with $z=c\alpha$.

\subsubsection{Finite-dimensional escort average}
Suppose now that \(M=\mathcal B(\mathcal H)\) with
\(\dim\mathcal H<\infty\), and identify \(\psi\) and \(\omega\) with their
density matrices \(\rho\) and \(\sigma\).  Under the hypotheses above, the
map
\[
  \beta\longmapsto\rho_\beta^{(c)}
\]
is continuous in trace norm on \(I_\alpha\), and hence is Bochner integrable
with respect to \(\mu_\alpha\).  Since \(\mu_\alpha\) is a probability
measure, its barycenter
\begin{equation}
  \overline\rho_{\alpha,c}
  :=
  \int_{I_\alpha}\rho_\beta^{(c)}\,
  \dd\mu_\alpha(\beta)
  \label{eq:escort-barycenter}
\end{equation}
is again a density matrix.  The generalized Donald identity for continuous ensembles
\cite[Lemma~4, Eq.~(16)]{Holevo:2004ots}, together with
Theorem~\ref{thm:fixed-ray-integral}, gives the exact decomposition
\begin{equation}
  D_{\alpha,z}(\rho\Vert\sigma)
  =
  D(\overline\rho_{\alpha,c}\Vert\sigma)
  +
  \int_{I_\alpha}
    D(\rho_\beta^{(c)}\Vert\overline\rho_{\alpha,c})\,
    \dd\mu_\alpha(\beta).
  \label{eq:donald-holevo-escort}
\end{equation}
This is easily seen by adding and subtracting
$\Tr(\rho_\beta^{(c)}\log\overline\rho_{\alpha,c})$ in the integrand and
using the barycenter definition. If $\overline\rho_{\alpha,c}$ is singular, let $v\in \ker \overline\rho_{\alpha,c}$. Then the continuous non-negative function
\begin{equation}
    \beta\longmapsto \langle{v, \rho_\beta^{(c)} v}\rangle
\end{equation}
has zero $\mu_\alpha$ integral. Since $\mu_\alpha$ has full support, the above function vanishes identically. Hence
\begin{equation}
    s(\rho_\beta^{(c)} ) \leq s(\overline\rho_{\alpha,c}),
\end{equation}
for every $\beta\in I_\alpha$. This support inclusion ensures that the add-and-subtract
calculation involving $\log\overline\rho_{\alpha,c}$ is legitimate, with
the logarithm taken on its support.



The second term in \eqref{eq:donald-holevo-escort} is the Holevo information of the
\(\mu_\alpha\)-weighted escort ensemble.  Denote it by
\begin{equation}
  \chi_{\alpha,c}
  :=
  \int_{I_\alpha}
    D(\rho_\beta^{(c)}
      \Vert\overline\rho_{\alpha,c})\,
    \dd\mu_\alpha(\beta).
  \label{eq:escort-holevo-information}
\end{equation}
Applying Pinsker's inequality pointwise to
\((\rho_\beta^{(c)},\overline\rho_{\alpha,c})\) and integrating gives
\begin{equation}
  \chi_{\alpha,c}
  \geq
  \frac12
  \int_{I_\alpha}
    \bigl\|
      \rho_\beta^{(c)}-\overline\rho_{\alpha,c}
    \bigr\|_1^2\,
    \dd\mu_\alpha(\beta).
  \label{eq:escort-holevo-pinsker}
\end{equation}
Equivalently, the Donald decomposition admits the stability refinement
\begin{equation}
  D_{\alpha,z}(\rho\Vert\sigma)
  \geq
  D(\overline\rho_{\alpha,c}\Vert\sigma)
  +
  \frac12
  \int_{I_\alpha}
    \bigl\|
      \rho_\beta^{(c)}-\overline\rho_{\alpha,c}
    \bigr\|_1^2\,
    \dd\mu_\alpha(\beta).
  \label{eq:donald-holevo-stability}
\end{equation}
Thus the Holevo term quantitatively controls the mean-square trace-norm
dispersion of the escort states around their barycenter.  It vanishes
precisely when
\(\rho_\beta^{(c)}=\overline\rho_{\alpha,c}\) for
\(\mu_\alpha\)-almost every \(\beta\).  By trace-norm continuity and the
full support of \(\mu_\alpha\) on \(I_\alpha\), this is equivalent to the
escort trajectory being constant on \(I_\alpha\).

\subsubsection{Algebraic rigidity of the escort Holevo term}
The almost-everywhere part of the preceding vanishing criterion does not
depend on finite dimensionality.  Let \(M\) be an arbitrary von Neumann
algebra and suppose that the escort path
\[
  \beta\longmapsto\psi_\beta^{(c)}\in M_*
\]
is strongly measurable as an \(M_*\)-valued map.  Since
\(\|\psi_\beta^{(c)}\|_{M_*}=1\), the path is Bochner integrable, and
\begin{equation}
  \overline\psi_{\alpha,c}
  :=
  \int_{I_\alpha}
    \psi_\beta^{(c)}\,\dd\mu_\alpha(\beta)
  \in M_*
  \label{eq:algebraic-escort-barycenter}
\end{equation}
is a normal state.  Moreover, the map
\[
  \varphi\longmapsto
  D(\varphi\Vert\overline\psi_{\alpha,c})
\]
is lower semicontinuous, and hence Borel measurable, for the
predual-norm topology.  Combined with the assumed strong
measurability of
\(\beta\mapsto\psi_\beta^{(c)}\), this shows that
\[
  \beta\longmapsto
  D\!\left(
    \psi_\beta^{(c)}
    \middle\Vert
    \overline\psi_{\alpha,c}
  \right)
\]
is measurable with respect to the \(\mu_\alpha\)-completion of
\(\mathcal B(I_\alpha)\), the Borel $\sigma$-algebra of $I_\alpha$.  Hence the extended nonnegative quantity
\begin{equation}
  \chi_{\alpha,c}
  :=
  \int_{I_\alpha}
    D\!\left(
      \psi_\beta^{(c)}
      \middle\Vert
      \overline\psi_{\alpha,c}
    \right)
    \,\dd\mu_\alpha(\beta)
  \label{eq:algebraic-escort-holevo}
\end{equation}
is well defined.  Pinsker's inequality for normal states on \(M\) \cite{Hiai1981SufficiencyKC} gives
\begin{equation}
  \chi_{\alpha,c}
  \geq
  \frac12
  \int_{I_\alpha}
    \bigl\|
      \psi_\beta^{(c)}-\overline\psi_{\alpha,c}
    \bigr\|_{M_*}^{\,2}
    \,\dd\mu_\alpha(\beta).
  \label{eq:algebraic-escort-holevo-pinsker}
\end{equation}
Consequently,
\begin{equation}
    \chi_{\alpha,c}=0
    \quad\Longleftrightarrow\quad
    \psi_\beta^{(c)}
    =
    \overline\psi_{\alpha,c}
    \quad\text{for \(\mu_\alpha\)-almost every
    \(\beta\in I_\alpha\)}.
  \label{eq:algebraic-escort-holevo-rigidity}
\end{equation}
Thus neither finite dimensionality nor continuity in the R\'enyi order is
needed for this almost-everywhere rigidity statement.  For this
equivalence alone, strong measurability may be replaced by the weaker
assumptions that the escort family admits a normal \(M_*\)-valued barycenter
and that the displayed relative-entropy profile is measurable.

This statement concerns the rigidity of the Holevo quantity whenever it is
defined.  Identifying \(\chi_{\alpha,c}\) as the exact remainder in a Donald
decomposition on an arbitrary von Neumann algebra would additionally require
a continuous-ensemble Donald identity in that setting.  For
\(0<\alpha<1\) without support inclusion and with \(q_c>0\), the same conclusion applies to the
compressed escort family; the separate support-boundary term is not part of
\(\chi_{\alpha,c}\).
 

\subsection{Binary testing and generalized cutoff rates}
\label{subsec:binary-testing}
Let $0\leq T\leq\one$ be a binary test on $M$ for distinguishing two
normal states $\psi,\omega$, and put
\begin{equation}
  p\defeq \psi(T),\qquad q\defeq \omega(T).
\end{equation}
Consider the normal unital completely-positive measurement channel
\begin{equation}
    \mathcal{M}_T:\mathbb{C}^2\to M,\qquad \mathcal{M}_T(a,b)=aT+b(\one-T).
\end{equation}
For $\alpha>1$ and $(\alpha,z)$ in the upper data-processing region of
\eqref{eq:dpi-region}, monotonicity under $\mathcal{M}_T$ and the classical binary
formula give, whenever $p,q>0$,
\begin{align}
  D_{\alpha,z}(\psi\Vert\omega)
  &\geq D_\alpha\bigl((p,1-p)\Vert(q,1-q)\bigr)\notag\\
  &\geq\frac{\alpha}{\alpha-1}\log p-\log q.
  \label{eq:one-shot-testing-derivation}
\end{align}
The second inequality follows by retaining the first nonnegative summand
in the classical R\'enyi moment.  For a nontrivial finite bound, assume
$D_{\alpha,z}(\psi\Vert\omega)<\infty$.  If $p=0$, the conclusion below
is automatic, while if $p>0$, measurement monotonicity forces $q>0$.
Rearranging gives
\begin{equation}
  -\log p
  \geq\frac{\alpha-1}{\alpha}
  \left[-\log q-D_{\alpha,z}(\psi\Vert\omega)\right].
  \label{eq:one-shot-testing-bound}
\end{equation}
This is valid on an arbitrary von Neumann algebra whenever $D_{\alpha,z}(\psi\Vert\omega)<\infty$,
$\alpha>1$ and $(\alpha,z)$ lies in the upper data-processing region
of \eqref{eq:dpi-region}. This bound states that the success
exponent is bounded below by a line of slope
\((\alpha-1)/\alpha\) in the type-II exponent, with horizontal
intercept \(D_{\alpha,z}(\psi\Vert\omega)\).  
The same one-copy measurement inequality underlies the proof of \cite[Lemma~IV.7]{Mosonyi:2014wsq}.   


For tensor powers, let
\(T_n\in M^{\bar\otimes n}\), \(0\leq T_n\leq\one\), and set
\[
  p_n\defeq\psi^{\otimes n}(T_n),
  \qquad
  q_n\defeq\omega^{\otimes n}(T_n).
\]
We use the standard tensor-product functoriality of Haagerup spaces.
Under the canonical realization associated with the spatial tensor product,
\begin{equation}
  h_{\psi^{\otimes n}}=h_\psi^{\otimes n},
  \qquad
  h_{\omega^{\otimes n}}=h_\omega^{\otimes n},
  \qquad
  s(\omega^{\otimes n})=s(\omega)^{\otimes n},
  \label{eq:tensor-haagerup-densities}
\end{equation}
and powers and closed products of positive elementary tensors are taken
componentwise.

From \cite[Proposition 10]{Kato:2023hlj}, one has
\begin{equation}
  D_{\alpha,z}
  \bigl(\psi^{\otimes n}\Vert\omega^{\otimes n}\bigr)
  =
  nD_{\alpha,z}(\psi\Vert\omega).
  \label{eq:tensor-alpha-z-additivity}
\end{equation}
Hence, if
$q_n\leq\e^{-nr}$ for every sufficiently large $n$, then
\begin{equation}
  \liminf_{n\to\infty}-\frac1n\log p_n
  \geq\frac{\alpha-1}{\alpha}
       \left[r-D_{\alpha,z}(\psi\Vert\omega)\right].
  \label{eq:iid-testing-converse}
\end{equation}
On a fixed ray this lower bound is the signed area
\begin{equation}
  \frac{\alpha-1}{\alpha}
  \left[r-D_{\alpha,c\alpha}(\psi\Vert\omega)\right]
  =\int_1^\alpha
   \frac{r-R_c(\beta)}{\beta^2}\,\dd\beta.
  \label{eq:generic-testing-area}
\end{equation}
For an upper-order fixed ray, the DPI condition is
\begin{equation}
  \max\!\left\{\frac12,1-\frac1\alpha\right\}
  \leq c\leq1.
  \label{eq:fixed-ray-dpi}
\end{equation}

The exact strong-converse exponent of asymmetric binary hypothesis testing
selects the sandwiched ray.  Define
\begin{equation}
  p_n^\star(r)
  \defeq \sup\left\{
  \psi^{\otimes n}(T_n):0\leq T_n\leq\one,\ 
  \omega^{\otimes n}(T_n)\leq\e^{-nr}
  \right\}.
  \label{eq:optimal-testing-success}
\end{equation}
Whenever the corresponding operational theorem applies, set
\begin{equation}
  \scal(r;\psi\Vert\omega)
  \defeq -\lim_{n\to\infty}\frac1n\log p_n^\star(r).
\end{equation}
For $r>D(\psi\Vert\omega)$, the known strong-converse formula for $\alpha>1$ is
\begin{equation}
  \scal(r;\psi\Vert\omega)
  =\sup_{\alpha\in\mathcal I_+}
  \frac{\alpha-1}{\alpha}
  \bigl(r-D_\alpha(\psi\Vert\omega)\bigr).
  \label{eq:known-strong-converse}
\end{equation}
Here
\begin{equation}
  \mathcal I_+
  \defeq \{\alpha>1:Q_\alpha(\psi\|\omega)<\infty\}.
  \label{eq:finite-upper-srd-domain}
\end{equation}
Equivalently, one may optimize over all $\alpha>1$ using extended values,
since an infinite divergence contributes $-\infty$. 
\eqref{eq:known-strong-converse} is known to hold in finite dimensions \cite{Mosonyi:2014wsq}, on an injective von Neumann algebra when
$D_{\alpha_0}(\psi\Vert\omega)<\infty$ for some $\alpha_0>1$
\cite{Hiai:2021slp}, and on an arbitrary von Neumann algebra under the
bounded-domination hypothesis $a\psi\leq\omega$ for some  $a>0$ \cite{Junge:2025kuz}.
% Under bounded
% domination, $\mathcal I_+=(1,\infty)$.

Combining \eqref{eq:known-strong-converse} with
Corollary~\ref{cor:srd-integral} gives the exact area formula
\begin{equation}
  \scal(r;\psi\Vert\omega)
  =\sup_{\alpha\in\mathcal I_+}
   \int_1^\alpha
   \frac{r-D(\psi_\beta^s\Vert\omega)}{\beta^2}\,\dd\beta .
  \label{eq:strong-converse-area}
\end{equation}
The profile
\[
  R_1(\beta)\defeq D(\psi_\beta^s\Vert\omega)
\]
is nondecreasing, so the right limit
\[
  R_1(1+)\defeq\lim_{\beta\downarrow1}R_1(\beta)
\]
exists in $[0,+\infty]$.  In each of the regimes considered above, the
relevant hypotheses ensure that
\[
  D_{\alpha_0}(\psi\Vert\omega)<+\infty
\]
for some $\alpha_0>1$: this is assumed explicitly in the injective case,
follows from bounded domination in the arbitrary-algebra case, and, in
finite dimensions, follows from $D(\psi\Vert\omega)<+\infty$.  The
right-endpoint theorem for the sandwiched R\'enyi divergence therefore
applies and gives
\begin{equation}
  \lim_{\alpha\downarrow1}
  D_\alpha(\psi\Vert\omega)
  =
  D(\psi\Vert\omega).
  \label{eq:upper-srd-endpoint-operational}
\end{equation}
See \cite{Berta:2016vnw,Jencova:2016tqz,Jencova:2017txf}.

On the other hand, for $1<\alpha\leq\alpha_0$,
Corollary~\ref{cor:srd-integral} and
\[
  \frac{\alpha}{\alpha-1}
  \int_1^\alpha\frac{\dd\beta}{\beta^2}=1
\]
show that $D_\alpha(\psi\Vert\omega)$ is a weighted average of
$R_1$ over $(1,\alpha)$.  Monotonicity of $R_1$ therefore gives
\[
  R_1(1+)
  \leq
  D_\alpha(\psi\Vert\omega)
  \leq
  R_1(\alpha).
\]
Letting $\alpha\downarrow1$ and using the definition of $R_1(1+)$ yields
\[
  \lim_{\alpha\downarrow1}
  D_\alpha(\psi\Vert\omega)
  =
  R_1(1+).
\]
Together with \eqref{eq:upper-srd-endpoint-operational}, this proves
\[
  \lim_{\beta\downarrow1}R_1(\beta)
  =
  D(\psi\Vert\omega).
\]

Hence, if $r>D(\psi\Vert\omega)$, the integrand in
\eqref{eq:strong-converse-area} is positive for all $\beta>1$ sufficiently
close to one.  It need not ever become negative, and the supremum may be
attained only at the boundary of $\mathcal I_+$, including in the limit
$\alpha\to\infty$.  If it is attained at a finite interior order
$\alpha_r$, the subgradient condition is
\[
  R_1(\alpha_r-)\leq r\leq R_1(\alpha_r+).
\]
At a continuity point this becomes
\begin{equation}
  r=R_1(\alpha_r)=D(\psi_{\alpha_r}^s\Vert\omega).
  \label{eq:strong-converse-level-crossing}
\end{equation}
We emphasize that the imported operational theorems
\cite{Mosonyi:2014wsq,Hiai:2021slp,Junge:2025kuz} supply the exponent, and the
new content of \eqref{eq:strong-converse-area} is its escort-area and
level-crossing form.

For $0<\kappa<1$, define the generalized cutoff rate by
\begin{equation}
  C_\kappa(\psi\Vert\omega)
  \defeq \inf\left\{r_0\in\mathbb R:
  \scal(r;\psi\Vert\omega)\geq\kappa(r-r_0)
  \text{ for every }r>0\right\}.
  \label{eq:cutoff-rate-definition}
\end{equation}
Under the bounded-domination hypothesis,
\cite{Junge:2025kuz} proves that these generalized cutoff rates
satisfy
$C_\kappa=D_{1/(1-\kappa)}(\psi\Vert\omega)$ for $0<\kappa<1$.
Consequently,
\begin{equation}
  C_\kappa(\psi\Vert\omega)
  =\frac1\kappa\int_0^\kappa
   D\!\left(\psi_{1/(1-u)}^s\Vert\omega\right)\,\dd u.
  \label{eq:cutoff-rate-average}
\end{equation}
This gives an exact operational meaning to the average of the upper-order
sandwiched escort profile, beyond optimization over the order.

\subsection{Finite-dimensional work and pair conversion}

We next record two applications whose established operational theorems are
finite-dimensional.  Let
\begin{equation}
  \tau_{\beta_{\rm th}}
  \defeq \frac{\e^{-\beta_{\rm th}H}}
          {\Tr\e^{-\beta_{\rm th}H}}
  \label{eq:gibbs-state}
\end{equation}
be a faithful Gibbs state at inverse temperature $\beta_{\rm th}>0$.

For the work-extraction application, we use the battery model of
\cite{Watanabe:2026pzw}.  For $m>1$, let $X_m$ be a two-level battery with
basis $\{|0\rangle,|1\rangle\}$, Hamiltonian gap
$\beta_{\rm th}^{-1}\log(m-1)$, and thermal state
\[
  \mu_m
  =\frac{m-1}{m}|0\rangle\!\langle0|
   +\frac1m|1\rangle\!\langle1|.
\]
The transition from $\rho\otimes\mu_m$ to the target pure battery state
$|1\rangle\!\langle1|$ is assigned work
$W=\beta_{\rm th}^{-1}\log m$. Define the corresponding
dimensionless work parameter by
\[
  w\defeq\beta_{\rm th}W=\log m,
  \qquad
  W=\beta_{\rm th}^{-1}w.
\]
Thus the rate $r$ below is the dimensionless work per copy.
Write
$H_{X_m}=\beta_{\rm th}^{-1}\log(m-1)|1\rangle\!\langle1|$.  For $w>0$,
let ${\rm TO}(H_w^{\rm in}\to H_w^{\rm out})$ denote the thermal operations\footnote{As in~\cite{Watanabe:2026pzw}, the class of thermal operations is
understood to include limits of ordinary thermal operations.}
with
\[
  H_w^{\rm in}=H\otimes\one+\one\otimes H_{X_{\e^w}},
  \qquad H_w^{\rm out}=H_{X_{\e^w}}.
\]
Define the optimal fidelity error by
\[
  \mathcal E_{\rm TO}(\rho;w)
  \defeq 1-\sup_{\Lambda\in{\rm TO}(H_w^{\rm in}\to H_w^{\rm out})}
  \langle1|\Lambda(\rho\otimes\mu_{\e^w})|1\rangle.
\]
Here \cite{Watanabe:2026pzw} uses squared fidelity, which reduces to the displayed
matrix element because the target is pure.  \cite[Lemma~S.15]{Watanabe:2026pzw} proves equivalence with a sharp work-storage model
for exact extraction, but whether this equivalence persists with nonzero
fidelity error is not known.  We therefore retain the thermal-battery
model throughout.

For $n$ copies, the system Hamiltonian is the additive Hamiltonian
$H^{\times n}\defeq\sum_{j=1}^n
\one^{\otimes(j-1)}\otimes H\otimes\one^{\otimes(n-j)}$.
When the argument of $\mathcal E_{\rm TO}$ is $\rho^{\otimes n}$, its
definition uses this system Hamiltonian in $H_w^{\rm in}$.
For $r>0$, define the reliability exponent
\begin{equation}
  B_{\rm TO}(\rho;r)
  \defeq \sup_{\substack{(w_n)\subset(0,\infty)\\
                         \liminf_{n\to\infty}w_n/n\geq r}}
  \liminf_{n\to\infty}
  \left[-\frac1n\log
  \mathcal E_{\rm TO}(\rho^{\otimes n};w_n)\right].
  \label{eq:work-reliability-definition}
\end{equation}
with the convention $-\log0=+\infty$.
\cite[Theorem~2]{Watanabe:2026pzw} gives
\begin{equation}
  B_{\rm TO}(\rho;r)
  =\sup_{0<\alpha<1}
   \frac{\alpha-1}{\alpha}
   \left(r-D_\alpha(\rho\Vert\tau_{\beta_{\rm th}})\right).
  \label{eq:work-reliability-renyi}
\end{equation}
The supremum includes $0<\alpha<1/2$, illustrating why the analytic range
of Corollary~\ref{cor:srd-integral} is useful even outside the SRD
data-processing range.  Substitution gives
\begin{equation}
  B_{\rm TO}(\rho;r)
  =\sup_{0<\alpha<1}
   \int_\alpha^1
   \frac{D(\rho_\beta^s\Vert\tau_{\beta_{\rm th}})-r}{\beta^2}
   \,\dd\beta .
  \label{eq:work-reliability-area}
\end{equation}
In finite dimensions the escort profile is continuous, hence a maximizer
at a finite interior order $\alpha_r$ satisfies
$D(\rho_{\alpha_r}^s\Vert\tau_{\beta_{\rm th}})=r$.
For Gibbs-preserving operations the reliability exponent is instead
governed by the Petz divergences
\(\overline D_\alpha=D_{\alpha,1}\).  Since the Petz line \(z=1\)
does not lie on a single ray \(z=c\alpha\), our fixed-ray
representation does not turn this optimization into a level crossing
against one escort profile.

There is also an operational setting in which $\alpha$ and $z$ vary
independently.  Let $\mathcal F$ be the finite-dimensional class of pairs
\begin{align}
  \rho&=\sum_i p_i\lvert i\rangle\!\langle i\rvert
      \otimes\lvert a_i\rangle\!\langle a_i\rvert, \qquad
  \sigma=\sum_i q_i\lvert i\rangle\!\langle i\rvert
      \otimes\lvert b_i\rangle\!\langle b_i\rvert
  \label{eq:flat-pair-class}
\end{align}
where $(p_i)_i$ and $(q_i)_i$ are probability distributions,
$|a_i\rangle$ is normalized when $p_i>0$ and is the zero vector when
$p_i=0$, and similarly $|b_i\rangle$ is normalized when $q_i>0$ and is
zero when $q_i=0$.  Membership in $\mathcal F$ includes the
\emph{some-overlap} condition: $\langle a_i,b_i\rangle\neq0$ for at least
one index $i$.  A source pair is \emph{nonparallel} when
$|\langle a_i,b_i\rangle|<1$ for every $i$.  In addition, the source pair
below is assumed to satisfy the condition in \cite[eq.~(14)]{Verhagen:2025anx}:
for some $i$, both vectors are nonzero and orthogonal.

Write $(\rho,\sigma)\succeq(\rho',\sigma')$ if a single channel maps
$\rho$ to $\rho'$ and $\sigma$ to $\sigma'$.  Put
$\mathsf P=(\rho,\sigma)$ and $\mathsf P'=(\rho',\sigma')$.  A sequence
$(m_n)_{n\geq1}\subset\mathbb N$ is \emph{achievable} if
\begin{equation}
  \mathsf P^{\otimes n}\succeq(\mathsf P')^{\otimes m_n}
  \quad\text{for every sufficiently large }n.
  \label{eq:achievable-pair-conversion-sequence}
\end{equation}
The exact large-sample conversion rate is
\begin{equation}
  r(\mathsf P\to\mathsf P')
  \defeq \sup\left\{
  \liminf_{n\to\infty}\frac{m_n}{n}:
  (m_n)\ \text{is achievable}\right\}.
  \label{eq:exact-large-sample-rate-definition}
\end{equation}
For source and target pairs in $\mathcal F$, with the source nonparallel
and satisfying the condition in \cite[eq.~(14)]{Verhagen:2025anx}, 
\cite[Theorem~5]{Verhagen:2025anx} gives
\begin{equation}
  r\bigl((\rho,\sigma)\to(\rho',\sigma')\bigr)
  =\inf_{\substack{0<\alpha<1\\
                    z>\max\{\alpha,1-\alpha\}}}
  \frac{D_{\alpha,z}(\rho\Vert\sigma)}
       {D_{\alpha,z}(\rho'\Vert\sigma')}.
  \label{eq:pair-conversion-rate}
\end{equation}
Here a quotient with zero denominator is interpreted as $+\infty$, since
the corresponding monotone places no restriction on the conversion rate.
Condition~\cite[eq.~(14)]{Verhagen:2025anx} makes every source numerator strictly positive, so no
$0/0$ ambiguity occurs.  
Indeed, under Condition~\cite[eq.~(14)]{Verhagen:2025anx} there exists an $i_0$ such that $p_{i_0},q_{i_0}>0$ with zero overlap of the vectors, so
\[
\begin{aligned}
  Q_{\alpha,z}(\rho\Vert\sigma)
  &=
  \sum_i p_i^\alpha q_i^{1-\alpha}
  |\langle a_i,b_i\rangle|^{2z} \\
  &\leq
  \sum_i p_i^\alpha q_i^{1-\alpha}
  -p_{i_0}^\alpha q_{i_0}^{1-\alpha} \\
  &<
  \sum_i p_i^\alpha q_i^{1-\alpha}
  \leq 1.
\end{aligned}
\]
The some-overlap condition gives $0< Q_{\alpha,z}(\rho\|\sigma).$ 
% \begin{equation}
%    0< Q_{\alpha,z}(\rho\|\sigma).
% \end{equation}
Hence 
\begin{equation}
    0<D_{\alpha,z}(\rho\|\sigma)<\infty.
\end{equation}
The conclusion is asserted here only under the hypotheses of the cited
theorem, and \eqref{eq:pair-conversion-rate} is not a conversion formula for
arbitrary state pairs.

Since pairs in \(\mathcal F\) need not have nested supports, applying our
escort representation to the divergences in
\eqref{eq:pair-conversion-rate} requires the finite-dimensional boundary
correction of Proposition~\ref{prop:fd-boundary-correction}.  If $P=s(\sigma)$ and
$c=z/\alpha$, direct compression gives
\[
  P\rho^{1/c}P
  =\sum_{i:q_i>0}p_i^{1/c}|\langle a_i,b_i\rangle|^2
    |i\rangle\!\langle i|\otimes|b_i\rangle\!\langle b_i|.
\]
Consequently, the boundary datum is
\begin{equation}
  q_{1-,c}
  =\sum_{i:q_i>0}p_i\lvert\langle a_i,b_i\rangle\rvert^{2c},
  \qquad c=\frac z\alpha.
  \label{eq:cq-boundary-datum}
\end{equation}
This number is strictly positive for every pair in $\mathcal F$, by the
some-overlap condition.  Thus \eqref{eq:pair-conversion-rate} is organized
by a family of escort profiles, one for each $c$, together with their
support-overlap costs; the boundary term vanishes whenever the relevant
source support is contained in the reference support.

\section{Lower-order sandwiched divergence}
\label{sec:proof-lower-srd}

We now prove the lower-order sandwiched integral formula directly in the
data-processing range $1/2\leq\alpha<1$.  This argument is worth retaining
even though the fixed-ray proof in Section~\ref{sec:proof-alpha-z}
subsumes it, since it uses only the Banach range of Araki--Masuda interpolation
\cite{Berta:2016vnw,Jencova:2017txf} and therefore does not depend
on quasi-Banach interpolation below $L^1$.

\subsection{Direct proof of integral representation}

For normal states $\psi, \omega$, assume first that $s(\psi)\leq s(\omega)$ and pass to
$s(\omega)Ms(\omega)$, so that $\omega$ is faithful and the restrictions of $\psi,\omega$ remain states.  For
$1/2\leq t<1$ let $A_t$ be the closed product in
\eqref{eq:lower-srd-factor}, that is,
\begin{equation}
    A_t\defeq x_t^s(\psi\|\omega),
\end{equation}
and set
\begin{equation}
  g(t)\defeq \log Q_t(\psi\|\omega)
       =\log\tr(A_t^t),
  \qquad g(1)\defeq 0.
  \label{eq:lower-proof-g}
\end{equation}

\begin{lemma}[Convexity of the lower log moment]
\label{lem:lower-log-convexity}
The function $g$ is finite and convex on $[1/2,1)$ and extends
continuously to $t=1$ with $g(1)=0$.
\end{lemma}

\begin{proof}
For $1\leq p\leq2$, put
\[
  a_p\defeq \frac1p-\frac12
\]
and define the closed Haagerup product
\begin{equation}
  B_p
  \defeq
  \overline{h_\omega^{a_p}h_\psi^{1/2}}
  \in L^p(M),
\end{equation}
together with
\begin{equation}
  N_p\defeq\|B_p\|_p.
  \label{eq:lower-weighted-norm}
\end{equation}
For $1\leq p<2$, if
\[
  q_p\defeq\frac{2p}{2-p},
\]
then $a_p=1/q_p$ and $1/p=1/q_p+1/2$.  Generalized
H\"older multiplication therefore gives
\[
  N_p
  \leq
  \|h_\omega^{a_p}\|_{q_p}
  \|h_\psi^{1/2}\|_2
  =
  \omega(\one)^{1/q_p}\psi(\one)^{1/2}
  =1.
\]
The same conclusion holds for $p=2$, since
$B_2=h_\psi^{1/2}$.

Under the canonical Haagerup standard-form identification, let
\[
  \xi_\psi=h_\psi^{1/2}\in L^2(M)
\]
denote the natural-cone vector representative of $\psi$.  The Haagerup
realization of the Araki--Masuda spaces
\cite{Berta:2016vnw,Jencova:2017txf,Kato:2023aro} gives
\begin{equation}
  N_p=\|\xi_\psi\|_{p,\omega}^{\rm AM},
  \qquad 1\leq p\leq2.
  \label{eq:lower-haagerup-am-identification}
\end{equation}
Since $\omega$ is faithful, $\|\cdot\|_{p,\omega}^{\rm AM}$ is a norm.
Moreover, $\psi\neq0$ implies $\xi_\psi\neq0$.  Consequently,
\[
  0<N_p\leq1,
  \qquad 1\leq p\leq2.
\]

Now let $p=2t$, where $1/2\leq t\leq1$.  Then
\[
  a_p=\frac{1-t}{2t},
\]
and associativity in the $*$-algebra of $\tau$-measurable operators
gives
\begin{equation}
  A_t
  =
  \overline{
    h_\omega^{a_p}h_\psi h_\omega^{a_p}}
  =
  B_pB_p^*.
  \label{eq:lower-factor-from-Bp}
\end{equation}
Let $B_p=v_p|B_p|$ be the polar decomposition of $B_p$.  Since
$p=2t$ and $|B_p|^p\in L^1(M)$, functional calculus and the tracial
property of $\tr$ for H\"older compatible elements gives
\begin{align}
  Q_t(\psi\|\omega)
  &=\tr(A_t^t) \notag\\
  &=\tr\bigl((B_pB_p^*)^{p/2}\bigr) \notag\\
  &=\tr\bigl(v_p|B_p|^p v_p^*\bigr) \notag\\
  &=\tr(|B_p|^p)
   =\|B_p\|_p^p,
\end{align}
using $v_p\in M$ and $v_p^*v_p=s(|B_p|)$. 
Thus
\begin{equation}
  Q_t(\psi\|\omega)=N_{2t}^{2t},
  \qquad \frac12\leq t\leq1.
  \label{eq:lower-moment-weighted-norm}
\end{equation}

Let $1\leq p_0,p_1\leq2$, $0\leq\lambda\leq1$, and set
\begin{equation}
  p=(1-\lambda)p_0+\lambda p_1,
  \qquad
  \theta=\frac{\lambda p_1}{p}.
\end{equation}
Then
\begin{equation}
  \frac1p=\frac{1-\theta}{p_0}
          +\frac{\theta}{p_1},
  \qquad
  p(1-\theta)=(1-\lambda)p_0,
  \qquad p\theta=\lambda p_1.
\end{equation}
The weighted Araki--Masuda interpolation inequality
\cite[Proposition~4]{Berta:2016vnw} yields
\begin{equation}
  N_p\leq N_{p_0}^{1-\theta}N_{p_1}^{\theta}.
\end{equation}
After raising this inequality to the power $p$, we obtain
\begin{equation}
  N_p^p
  \leq (N_{p_0}^{p_0})^{1-\lambda}
       (N_{p_1}^{p_1})^\lambda.
\end{equation}
Together with \eqref{eq:lower-moment-weighted-norm}, this proves convexity
of $g$.  Finiteness was already obtained from H\"older's inequality, and
$N_2=\|h_\psi^{1/2}\|_2=1$ gives the assigned value $g(1)=0$.

It remains to rule out a jump at $t=1$, since convexity alone does not do so.
Let $\overline Q_t(\psi\|\omega)$ denote the Petz moment.  The
Araki--Lieb--Thirring inequality in standard form
\cite[Theorem~12]{Berta:2016vnw} gives, for
$1/2\leq t<1$,
\begin{equation}
  \log\overline Q_t(\psi\|\omega)
  \leq g(t)\leq0.
  \label{eq:petz-sandwiched-endpoint-squeeze}
\end{equation}
If $\Delta_{\omega,\psi}$ is the relative modular operator, then
\begin{equation}
  \overline Q_t(\psi\|\omega)
  =\left\langle\xi_\psi,
  \Delta_{\omega,\psi}^{1-t}\xi_\psi\right\rangle.
  \label{eq:petz-standard-form-moment}
\end{equation}
Let $E_{\omega,\psi}$ denote the projection-valued spectral measure of
$\Delta_{\omega,\psi}$, and define
\[
  \mu_\psi(B)
  \defeq
  \left\langle\xi_\psi,
  E_{\omega,\psi}(B)\xi_\psi\right\rangle,
\]
with Borel $B\subseteq[0,\infty)$.
Denote by $J$ the modular conjugation of the Haagerup standard form.  The
support formula for the relative modular operator
\cite[Theorem~2.4(1)]{Araki:1977zsq} is
\begin{equation}
  \one-E_{\omega,\psi}(\{0\})
  =
  s(\omega)J s(\psi)J.
  \label{eq:relative-modular-support}
\end{equation}
Indeed, $s(\omega)$ is the $M$-support of $\xi_\omega$, while
$Js(\psi)J$ is the $M'$-support of $\xi_\psi$.  Since
$\xi_\psi$ belongs to the natural cone, $J\xi_\psi=\xi_\psi$, and
$s(\psi)\xi_\psi=\xi_\psi$.  Hence
\[
  Js(\psi)J\xi_\psi=\xi_\psi.
\]
It follows from \eqref{eq:relative-modular-support} that
\begin{align}
  E_{\omega,\psi}(\{0\})\xi_\psi
  &=
  \bigl(\one-s(\omega)Js(\psi)J\bigr)\xi_\psi \notag\\
  &=
  \bigl(\one-s(\omega)\bigr)\xi_\psi.
  \label{eq:relative-modular-kernel-vector}
\end{align}
Consequently,
\begin{align}
  \mu_\psi(\{0\})
  &=
  \left\|E_{\omega,\psi}(\{0\})\xi_\psi\right\|^2 \notag\\
  &=
  \|(\one-s(\omega))\xi_\psi\|^2 \notag\\
  &=
  \psi(\one-s(\omega))
   =1-\psi(s(\omega)).
  \label{eq:relative-modular-zero-mass}
\end{align}
At the beginning of the proof we passed to the support corner
$s(\omega)Ms(\omega)$, in which $s(\omega)=\one$.  Hence
\[
  \mu_\psi(\{0\})=0.
\]
Since $\psi$ is a state, $\mu_\psi([0,\infty))=\|\xi_\psi\|^2=1$;
therefore $\mu_\psi$ is a probability measure concentrated on
$(0,\infty)$.

Fix
$t_0\in[1/2,1)$.  For $t_0\leq t<1$ and $\lambda>0$,
\begin{equation}
  \lambda^{1-t}\leq1+\lambda^{1-t_0}.
\end{equation}
The right-hand side is integrable because
$\overline Q_{t_0}(\psi\|\omega)<\infty$ from \eqref{eq:petz-sandwiched-endpoint-squeeze}.  Dominated convergence in
\eqref{eq:petz-standard-form-moment} therefore yields
$\overline Q_t(\psi\|\omega)\to1$ as $t\uparrow1$.  The squeeze
\eqref{eq:petz-sandwiched-endpoint-squeeze} proves
$g(t)\to0=g(1)$ and supplies the required endpoint continuity, even when
$D(\psi\Vert\omega)=+\infty$.
\end{proof}

For $\beta\in(1/2,1)$ define the unnormalized escort functional
$\nu_\beta\in M_*^+$ by
\begin{equation}
  h_{\nu_\beta}\defeq A_\beta^\beta.
  \label{eq:lower-unnormalized-escort}
\end{equation}
Then
\begin{equation}
  \nu_\beta(\one)=Q_\beta(\psi\|\omega),
  \qquad
  \nu_\beta=Q_\beta(\psi\|\omega)\psi_\beta^s.
  \label{eq:lower-escort-scaling}
\end{equation}

\begin{lemma}[Moving characteristic]
\label{lem:lower-moving-characteristic}
For $\beta\in(1/2,1)$ and
$r\in(1/(2\beta),1)$,
\begin{align}
  Q_{r,r\beta}(\nu_\beta\|\omega)
  &=Q_{r\beta}(\psi\|\omega),
  \label{eq:lower-moving-moment}\\
  D_{r,r\beta}(\nu_\beta\Vert\omega)
  &=\frac{g(r\beta)-g(\beta)}{r-1}.
  \label{eq:lower-moving-divergence}
\end{align}
\end{lemma}

\begin{proof}
Since $r<1$, the lower branch of the $\alpha$--$z$ definition gives
\begin{align}
  Q_{r,r\beta}(\nu_\beta\|\omega)
  &=\tr\!\left[
  \left(
  h_\omega^{\frac{1-r}{2r\beta}}
  h_{\nu_\beta}^{1/\beta}
  h_\omega^{\frac{1-r}{2r\beta}}
  \right)^{r\beta}\right]\notag\\
  &=\tr\!\left[
  \left(
  h_\omega^{\frac{1-r}{2r\beta}}
  A_\beta
  h_\omega^{\frac{1-r}{2r\beta}}
  \right)^{r\beta}\right].
  \label{eq:lower-moving-computation}
\end{align}
All products displayed here are closed products.  Associativity in the $*$-algebra of $\tau$ measurable operators and
\begin{equation}
  \frac{1-r}{2r\beta}+\frac{1-\beta}{2\beta}
  =\frac{1-r\beta}{2r\beta},
\end{equation}
identify the operator in parentheses with $A_{r\beta}$.  This proves
\eqref{eq:lower-moving-moment}.  Since
$\nu_\beta(\one)=\exp g(\beta)$, the normalized convention
\eqref{eq:normalized-alpha-z-positive} gives
\eqref{eq:lower-moving-divergence}.
\end{proof}

\begin{lemma}[Endpoint of the moving characteristic]
\label{lem:lower-moving-endpoint}
For every nonzero $\nu\in M_*^+$ and every $\beta>0$,
\begin{equation}
  \lim_{r\uparrow1}D_{r,r\beta}(\nu\Vert\omega)
  =D_1(\nu\Vert\omega).
  \label{eq:lower-moving-endpoint}
\end{equation}
\end{lemma}

\begin{proof}
For $r$ sufficiently close to one,
$\beta/2<r\beta<\beta$.  For fixed $0<r<1$, the lower-order divergence
is nondecreasing in its second parameter
\cite[Theorem~5.1]{Hiai:2024qve}; hence
\begin{equation}
  D_{r,\beta/2}(\nu\Vert\omega)
  \leq D_{r,r\beta}(\nu\Vert\omega)
  \leq D_{r,\beta}(\nu\Vert\omega).
  \label{eq:lower-moving-squeeze}
\end{equation}
For each fixed $z>0$, both outer terms converge to normalized Araki
relative entropy as $r\uparrow1$
\cite[Theorem~6.7]{Hiai:2024qve}.  The result follows by the squeeze theorem.
\end{proof}

We can now prove the integral formula.  Combining
Lemmas~\ref{lem:lower-moving-characteristic} and
\ref{lem:lower-moving-endpoint} gives, for
$\beta\in(1/2,1)$,
\begin{equation}
  D_1(\nu_\beta\Vert\omega)
  =\lim_{r\uparrow1}
   \frac{g(r\beta)-g(\beta)}{r-1}
  =\beta g'_-(\beta).
  \label{eq:lower-left-derivative}
\end{equation}
Since a finite convex function has finite one-sided derivatives at interior points, $ D_1(\nu_\beta\Vert\omega)<\infty$.
By \eqref{eq:lower-escort-scaling} and the relative-entropy scaling law
\eqref{eq:relative-entropy-scaling},
\begin{equation}
  D_1(\nu_\beta\Vert\omega)
  =D(\psi_\beta^s\Vert\omega)+g(\beta).
  \label{eq:lower-relative-entropy-scaling}
\end{equation}
Thus
\begin{equation}
  D(\psi_\beta^s\Vert\omega)
  =\beta g'_-(\beta)-g(\beta).
  \label{eq:lower-one-sided-refined}
\end{equation}
In particular, at every differentiability point of $g$, this is the
algebraic version of \eqref{eq:fd-refined-srd}.

The left derivative of a convex function is Borel measurable.  Moreover,
Lemma~\ref{lem:lower-log-convexity} and the convex endpoint fundamental
theorem (see Lemma~\ref{lem:convex-endpoint-ftc}), applied to the reflected function $u\mapsto g(1-u)$, imply that
$g$ is absolutely continuous on $[\alpha,1]$ whenever
$1/2<\alpha<1$.  Hence, for almost every $\beta$,
\begin{equation}
  \frac{\dd}{\dd\beta}\left(\frac{g(\beta)}\beta\right)
  =\frac{\beta g'(\beta)-g(\beta)}{\beta^2}
  =\frac{D(\psi_\beta^s\Vert\omega)}{\beta^2}.
\end{equation}
Integrating from $\alpha$ to $1$, using $g(1)=0$ and
$g(\alpha)=(\alpha-1)D_\alpha(\psi\Vert\omega)$, proves
\eqref{eq:main-srd-integral} for $1/2<\alpha<1$.  To reach
$\alpha=1/2$, let $\alpha_n\downarrow1/2$.  The established continuity of
the Araki--Masuda divergence at its lower endpoint
\cite[Lemma~8]{Berta:2016vnw} gives
$D_{\alpha_n}(\psi\Vert\omega)\to D_{1/2}(\psi\Vert\omega)$, while
monotone convergence applies to the nonnegative integrand on the expanding
intervals $[\alpha_n,1]$.  Since
$\alpha_n/(1-\alpha_n)\to1$, the same identity follows at
$\alpha=1/2$.  The value of the integrand at $\beta=1/2$ or $1$ is
immaterial to the Lebesgue integral.

For $0<\alpha<1/2$, the same identity follows by setting $c=1$ in the
all-$p$ proof of Theorem~\ref{thm:fixed-ray-integral}, given in Section~\ref{sec:proof-alpha-z}.

\subsection{Proof of the sandwiched support-boundary formula}

We end this section by proving
Corollary~\ref{cor:srd-support-boundary}.
Let $e=s(\omega)$ and $q=\psi(e)$, and suppose first that $q>0$.
Under the canonical identification
\begin{equation}
  L^1(eMe)\cong eL^1(M)e,
  \label{eq:corner-L1-identification}
\end{equation}
the canonical $L^1$-functionals and the corresponding
$L^1$--$L^\infty$ pairings are compatible. 
Denote the canonical functionals
temporarily by $\tr_e$ and $\tr_M$. The compatibility states that for $h\in eL^1(M)e, x\in eMe$,
\begin{equation}
    \tr_e(hx) = \tr_M(hx).
\end{equation}

Since the support of $h_\omega$ is $e$, one has
\[
  e h_\omega e=h_\omega.
\]
Thus, for every $x\in eMe$,
\[
  \tr_e(h_\omega x)
  =
  \tr_M(h_\omega x)
  =
  \omega(x)
  =
  \omega_e(x).
\]
Uniqueness of Haagerup $L^1$-densities therefore gives
\[
  h_{\omega_e}=h_\omega.
\]

Similarly, for every $x\in eMe$, the tracial property of the canonical
functional for the $L^1$--$L^\infty$ pairing gives
\begin{align}
  \tr_e\!\left(q^{-1}(e h_\psi e)x\right)
  &=
  q^{-1}\tr_M(eh_\psi e x) \notag\\
  &=
  q^{-1}\tr_M(h_\psi x) \notag\\
  &=
  q^{-1}\psi(x)
  =
  \widehat\psi(x).
  \label{eq:corner-density-pairing}
\end{align}
Again by uniqueness of the Haagerup density,
\[
  h_{\widehat\psi}=q^{-1}e h_\psi e.
\]
Consequently, under the canonical corner identification,
\begin{equation}
  h_{\omega_e}=h_\omega,
  \qquad
  e h_\psi e=q h_{\widehat\psi}.
  \label{eq:corner-density-scaling}
\end{equation}
We henceforth suppress the subscripts on the canonical functional $\tr$.

For $0<\beta<1$, the powers of $h_\omega$ are supported on $e$, so
\begin{align}
  x_\beta^s(\psi\|\omega)
  &=\overline{
  h_\omega^{\frac{1-\beta}{2\beta}}
  e h_\psi e
  h_\omega^{\frac{1-\beta}{2\beta}}}\notag\\
  &=q x_\beta^s(\widehat\psi\|\omega_e).
  \label{eq:corner-srd-factor}
\end{align}
It follows that
\begin{equation}
  Q_\beta(\psi\|\omega)
  =q^\beta Q_\beta(\widehat\psi\|\omega_e),
  \label{eq:corner-srd-moment}
\end{equation}
while the factor $q^\beta$ cancels in the normalized escort.  At
$\beta=\alpha$,
\begin{equation}
  D_\alpha(\psi\Vert\omega)
  =D_\alpha(\widehat\psi\Vert\omega_e)
   -\frac{\alpha}{1-\alpha}\log q.
  \label{eq:corner-srd-divergence}
\end{equation}
Applying the supported formula to $(\widehat\psi,\omega_e)$ proves
\eqref{eq:srd-support-boundary}.  For $\alpha\geq1/2$ this uses the direct
argument above; for $0<\alpha<1/2$ it uses
Theorem~\ref{thm:fixed-ray-integral}.  
If $q=0$, then $eh_\psi e\in L^1(M)_+$ and
\[
  \tr(eh_\psi e)=\psi(e)=q=0.
\]
Faithfulness of $\tr$ on $L^1(M)_+$ therefore gives
$eh_\psi e=0$.  Hence the lower sandwiched factor and moment vanish,
and the divergence is $+\infty$.
Finally,
$q=1$ is equivalent to $s(\psi)\leq e$, proving the last assertion of 
Corollary~\ref{cor:srd-support-boundary}.

\section{Proof of the fixed-ray \texorpdfstring{$\alpha$--$z$}{alpha-z}
representation}
\label{sec:proof-alpha-z}

We now prove Theorem~\ref{thm:fixed-ray-integral}, Proposition~\ref{prop:general-fixed-ray-boundary}, and Corollaries~\ref{cor:srd-integral},~\ref{cor:srd-support-boundary},~\ref{cor:raywise-monotonicity}.  We prove Theorem~\ref{thm:fixed-ray-integral} by realizing the fixed-ray factors as a compatible weighted Haagerup path.  An exact quotient identity and a varying-$z$ endpoint limit identify the left derivative of its logarithmic norm with the Araki relative entropy of the normalized power, while convexity supplies the absolute continuity needed to integrate this identity. A branch-dependent choice of anchor then identifies the path with the fixed-ray factors, yielding the theorem and Corollary~\ref{cor:raywise-monotonicity}.  Proposition~\ref{prop:general-fixed-ray-boundary} then follows by passing to the support corner and tracking the resulting normalization through the moments and escort states.


\subsection{A compatible weighted Haagerup scale}
\label{subsec:compatible-scale}

The following lemma makes the quasi-Banach input precise. 
All identities between closed Haagerup products below use the standard
associativity and generalized H\"older theorems for measurable operators;
see \cite{Terp:1981lp,Kato:2023aro}.

\begin{lemma}[Compatible weighted scale]
\label{lem:compatible-weighted-scale}
Let $\omega$ be a faithful normal state, let $r>0$, and let
$0\ne Y_r\in L^r(M)_+$.  For $0<p\leq r$, set
\begin{equation}
  d(p,r)\defeq \frac12\left(\frac1p-\frac1r\right)
  \label{eq:weighted-exponent}
\end{equation}
and put
\[
  T_{p,r}\defeq
  \overline{Y_r^{1/2}h_\omega^{d(p,r)}}\in L^{2p}(M).
\]
Define
the positive closed product
\begin{equation}
  Y_p
  \defeq T_{p,r}^*T_{p,r}
  =\overline{h_\omega^{d(p,r)}Y_rh_\omega^{d(p,r)}}
  \in L^p(M)_+.
  \label{eq:weighted-path}
\end{equation}
Then $0\ne Y_p\in L^p(M)_+$ and:
\begin{enumerate}[label=\textup{(\roman*)}]
\item for $0<p\leq q\leq r$,
\begin{equation}
  \overline{h_\omega^{d(p,q)}Y_qh_\omega^{d(p,q)}}=Y_p,
  \qquad
  \|Y_p\|_p\leq\|Y_q\|_q;
  \label{eq:weighted-compatibility}
\end{equation}
\item if
\begin{equation*}
  \frac1{p_\theta}
  =\frac{1-\theta}{p_0}+\frac{\theta}{p_1},
  \qquad 0<p_0<p_1\leq r,\qquad 0<\theta<1,
\end{equation*}
then
\begin{equation}
  \|Y_{p_\theta}\|_{p_\theta}
  \leq\|Y_{p_0}\|_{p_0}^{1-\theta}
       \|Y_{p_1}\|_{p_1}^{\theta};
  \label{eq:weighted-log-convexity}
\end{equation}
\item the anchoring endpoint is continuous:
\begin{equation}
  \lim_{p\uparrow r}\|Y_p\|_p=\|Y_r\|_r.
  \label{eq:weighted-endpoint-continuity}
\end{equation}
Consequently,
\begin{equation}
  u\longmapsto B(u)\defeq \log\|Y_{1/u}\|_{1/u}
  \label{eq:weighted-convex-B}
\end{equation}
is finite and convex for $u>1/r$ and has a finite continuous extension to
$u=1/r$.
\end{enumerate}
\end{lemma}

\begin{proof}
The factorization in
\eqref{eq:weighted-path} follows from closed multiplication and H\"older,
because
\begin{equation}
  \frac1{2p}=\frac1{2r}+d(p,r).
\end{equation}
To verify nonvanishing, the case $p=r$ is immediate, so suppose
$p<r$ and put
\[
  e_n=\one_{[1/n,n]}(h_\omega).
\]
Since $\omega$ is faithful, $e_n\uparrow\one$. If $T_{p,r}=0$, then
right multiplication by the bounded operator
$h_\omega^{-d(p,r)}e_n$ gives
\[
  Y_r^{1/2}e_n=0 .
\]
Let $k_r=\one_{\{0\}}(Y_r)$ be the kernel projection of $Y_r$.
The preceding equality implies $e_n\leq k_r$ for every $n$. Hence
\[
  \one=\bigvee_n e_n\leq k_r,
\]
so $k_r=\one$ and therefore $Y_r=0$. This contradicts the assumption
that $Y_r\neq0$. Thus $T_{p,r}\neq0$, and the factorization in
\eqref{eq:weighted-path} implies that $Y_p\neq0$.

Denote the modular automorphism group of $\omega$ by $\sigma^\omega$.
Let $\mathcal A_\omega\subset M$ be the subalgebra of 
$\sigma^\omega$-entire analytic elements and, for $a\in\mathcal A_\omega$, define the
closed product
\begin{equation}
  i_p^\omega(a)
  \defeq \overline{h_\omega^{1/(2p)}a h_\omega^{1/(2p)}}.
  \label{eq:compatible-core-embedding}
\end{equation}
By \cite[Lemma~1.1]{Junge:2003nbr}, the image
$i_p^\omega(\mathcal A_\omega)$ is dense in $L^p(M)$ for every $0<p<\infty$.
For $x=i_q^\omega(a)$ set
\begin{equation}
  J_{p,q}^\omega\bigl(i_q^\omega(a)\bigr)
  \defeq i_p^\omega(a),
  \qquad 0<p\leq q.
  \label{eq:downward-map}
\end{equation}
Since
\begin{equation}
  i_p^\omega(a)
  =\overline{h_\omega^{d(p,q)}
             i_q^\omega(a)
             h_\omega^{d(p,q)}},
\end{equation}
generalized H\"older multiplication gives, with $x=i_q^\omega(a)$ and the evident convention when $d(p,q)=0$,
\begin{align}
  \|J_{p,q}^\omega(x)\|_p
  &\leq
  \|h_\omega^{d(p,q)}\|_{1/d(p,q)}^2\|x\|_q\notag\\
  &=\tr(h_\omega)^{2d(p,q)}\|x\|_q
  =\|x\|_q.
  \label{eq:downward-contraction}
\end{align}
\cite[Proposition~2.6]{Gu:2019jfc} also records this compatible core
contraction.  For arbitrary $x\in L^q(M)$, generalized H\"older
multiplication further shows that the closed product
\begin{equation}
  \mathsf M_{p,q}(x)
  \defeq\overline{h_\omega^{d(p,q)}x h_\omega^{d(p,q)}}
  \label{eq:downward-closed-product}
\end{equation}
belongs to $L^p(M)$ and satisfies
$\|\mathsf M_{p,q}(x)\|_p\leq\|x\|_q$.  Since it agrees with
$J_{p,q}^\omega$ on the dense core, it is precisely the unique continuous
extension of that map.  We continue to denote the extension by
$J_{p,q}^\omega$.

For \(a\in\mathcal A_\omega\), the defining relation gives
\[
  \bigl(J^\omega_{p,q} J^\omega_{q,r}\bigr)
  \bigl(i_r^\omega(a)\bigr)
  =
  J^\omega_{p,q}\bigl(i_q^\omega(a)\bigr)
  =
  i_p^\omega(a)
  =
  J^\omega_{p,r}\bigl(i_r^\omega(a)\bigr).
\]
Since \(i_r^\omega(\mathcal A_\omega)\) is dense in \(L^r(M)\) and both
\(J^\omega_{p,q} J^\omega_{q,r}\) and \(J^\omega_{p,r}\) are
continuous maps from \(L^r(M)\) to \(L^p(M)\), the identity extends to
all of \(L^r(M)\):
\begin{equation}
  J^\omega_{p,q} J^\omega_{q,r}
  =
  J^\omega_{p,r}.
  \label{eq:downward-composition}
\end{equation}
By \eqref{eq:weighted-path} and the closed-product realization of
\(J^\omega_{s,r}\), we have
\[
  Y_s=J^\omega_{s,r}(Y_r),
  \qquad 0<s\le r.
\]
Consequently, for \(0<p\le q\le r\), the composition identity gives
\[
  J^\omega_{p,q}(Y_q)
  =
  J^\omega_{p,q}J^\omega_{q,r}(Y_r)
  =
  J^\omega_{p,r}(Y_r)
  =
  Y_p.
\]
In terms of closed products, this is
\[
  h_\omega^{d(p,q)}Y_qh_\omega^{d(p,q)}=Y_p.
\]
Finally, the contractivity of \(J^\omega_{p,q}\) yields
\[
  \|Y_p\|_p
  =
  \|J^\omega_{p,q}(Y_q)\|_p
  \leq
  \|Y_q\|_q.
\]
This proves 
\eqref{eq:weighted-compatibility}.  

We now use \cite[Theorem~4.1 and Corollary~4.3]{Gu:2019jfc} which provide the all-$p$ interpolation inequality
\begin{equation}
  \|i_{p_\theta}^\omega(a)\|_{p_\theta}
  \leq
  \|i_{p_0}^\omega(a)\|_{p_0}^{1-\theta}
  \|i_{p_1}^\omega(a)\|_{p_1}^{\theta},
  \label{eq:core-all-p-interpolation}
\end{equation}
proved by quasi-Banach complex interpolation in the crossed product.

Choose $a_n\in\mathcal A_\omega$ with
$i_r^\omega(a_n)\to Y_r$ in $L^r(M)$.  Contractivity gives
\begin{equation}
  i_p^\omega(a_n)
  =J_{p,r}^\omega(i_r^\omega(a_n))
  \longrightarrow Y_p
  \quad\text{in }L^p(M).
  \label{eq:core-approximation-all-p}
\end{equation}
For $0<p<1$, the standard Haagerup $L^p$ quasi-norm satisfies the
$p$-triangle inequality
\begin{equation}
  \|u+v\|_p^p
  \leq
  \|u\|_p^p+\|v\|_p^p,
  \qquad u,v\in L^p(M),
  \label{eq:haagerup-p-triangle}
\end{equation}
see \cite[Theorem~4.9(iii)]{Fack:1986ezo}.  Applying
\eqref{eq:haagerup-p-triangle} to
$x_n=x+(x_n-x)$ and then to $x=x_n+(x-x_n)$ gives
\begin{equation}
  \left|\|x_n\|_p^p-\|x\|_p^p\right|
  \leq
  \|x_n-x\|_p^p.
  \label{eq:p-power-continuity}
\end{equation}
Consequently, $x_n\to x$ in $L^p(M)$ implies
$\|x_n\|_p^p\to\|x\|_p^p$, and hence
$\|x_n\|_p\to\|x\|_p$.
Passing to the limit in \eqref{eq:core-all-p-interpolation} therefore
proves \eqref{eq:weighted-log-convexity} for the arbitrary anchor $Y_r$.

It remains to prove \eqref{eq:weighted-endpoint-continuity}.  For a fixed
$a\in\mathcal A_\omega$, choose $s>r$.  Contractivity gives
\begin{equation}
  \|i_p^\omega(a)\|_p\leq\|i_r^\omega(a)\|_r,
  \qquad p<r.
  \label{eq:core-endpoint-upper}
\end{equation}
Interpolate to $r$ between $p$ and $s$ to obtain
\begin{equation}
  \|i_r^\omega(a)\|_r
  \leq\|i_p^\omega(a)\|_p^{1-\theta(p)}
       \|i_s^\omega(a)\|_s^{\theta(p)},
  \qquad \theta(p)\longrightarrow0
  \quad(p\uparrow r),
  \label{eq:core-endpoint-lower}
\end{equation}
with $\theta(p)$ defined to satisfy $r^{-1}=(1-\theta(p))p^{-1}+\theta(p)s^{-1}$. 
Combining \eqref{eq:core-endpoint-upper} and \eqref{eq:core-endpoint-lower},
\begin{equation}
  \lim_{p\uparrow r}\|i_p^\omega(a)\|_p
  =\|i_r^\omega(a)\|_r.
  \label{eq:core-endpoint-continuity}
\end{equation}

Set
$\delta_n=\|Y_r-i_r^\omega(a_n)\|_r\to0$.
By contractivity,
\begin{equation}
  \|Y_p-i_p^\omega(a_n)\|_p\leq\delta_n
  \quad\text{uniformly for }p\leq r.
  \label{eq:endpoint-uniform-error}
\end{equation}
If $r>1$, take $p$ close enough to $r$ that $p\geq1$ and use the reverse
triangle inequality together with
\eqref{eq:core-endpoint-continuity}; this gives
\begin{equation}
  \limsup_{p\uparrow r}
  \left|\|Y_p\|_p-\|Y_r\|_r\right|
  \leq2\delta_n.
\end{equation}
Letting $n\to\infty$ proves the claim.

If $0<r\leq1$, take $p\in[r/2,r]$ and $0<\delta_n\leq1$.  By
\eqref{eq:p-power-continuity} and
\eqref{eq:endpoint-uniform-error},
\begin{equation}
  \left|\|Y_p\|_p^p-
  \|i_p^\omega(a_n)\|_p^p\right|
  \leq\delta_n^p\leq\delta_n^{r/2}.
\end{equation}
The analogous error at $p=r$ is at most $\delta_n^r$.  The ``reverse
triangle inequality'' together with
\eqref{eq:core-endpoint-continuity} now gives
\begin{equation}
  \limsup_{p\uparrow r}
  \left|\|Y_p\|_p^p-\|Y_r\|_r^r\right|
  \leq\delta_n^{r/2}+\delta_n^r.
\end{equation}
Letting $n\to\infty$ proves convergence of the $p$th powers.  Since
$p\to r>0$ and $Y_r\ne0$, taking the varying root gives
$\|Y_p\|_p\to\|Y_r\|_r$, proving endpoint continuity also in the
quasi-Banach case.  Finally, \eqref{eq:weighted-log-convexity} is exactly convexity of
$B$ in the reciprocal variable, and
\eqref{eq:weighted-endpoint-continuity} supplies its endpoint value.
\end{proof}

\begin{remark}[The quasi-Banach input]
\label{rem:quasi-banach-input}
The all-$p$ inequality \eqref{eq:core-all-p-interpolation} is a substantive
input from the preprint \cite{Gu:2019jfc}.  Without the quasi-Banach input, the
same proof is covered by ordinary Banach interpolation
\cite{KOSAKI198429,Terp:1981lp} whenever
the entire $p$ interval lies in $[1,\infty)$.  Concretely, this includes
the lower branch when $z\geq1$ and the upper branch when $z\geq\alpha$.
The separate argument in Section~\ref{sec:proof-lower-srd} also proves the
sandwiched formula for $1/2\leq\alpha<1$ without the all-$p$ input.  The
remaining unrestricted sub-$L^1$ range is exactly where
\cite{Gu:2019jfc} is used.
\end{remark}

We shall also use the following elementary convex-analysis fact.

\begin{lemma}[Convex endpoint fundamental theorem]
\label{lem:convex-endpoint-ftc}
Let $f$ be finite and convex on $(a,b)$ and suppose that it has a finite
continuous extension to $a$.  Then its extension is absolutely continuous
on every compact interval $[a,b']\subset[a,b)$, and
\begin{equation}
  f(b')-f(a)=\int_a^{b'}f'_+(u)\,\dd u
             =\int_a^{b'}f'_-(u)\,\dd u.
  \label{eq:convex-endpoint-ftc}
\end{equation}
The two one-sided derivatives agree at all but at most countably many points.
\end{lemma}

\begin{proof}
On every interior compact interval this is the standard fundamental
theorem for convex functions.  The one-sided derivatives are monotone.
Applying the interior formula on $[a+\varepsilon,b']$ and sending
$\varepsilon\downarrow0$, continuity at $a$ shows that the improper
integral of either derivative converges to $f(b')-f(a)$.  Its positive and
negative parts cannot both be infinite because a monotone function changes
sign at most once; since the limiting difference is finite, neither part
can be infinite.  Thus the derivative is integrable, which proves
absolute continuity and \eqref{eq:convex-endpoint-ftc}.
\end{proof}

The compatible weighted scale has the following consequence 
independent of the fixed-ray construction: along any compatible weighted $L^p$-path, the logarithmic $L^p$-norm is a potential whose infinitesimal slope is proportional to the Araki relative entropy of the normalized $p$-th power.

\begin{proposition}[Anchored entropy identity]
\label{prop:anchored-entropy-identity}
Let $\omega$ be a faithful normal state on $M$, let $r>0$, and let
$0\ne Y_r\in L^r(M)_+$.  For $0<p\leq r$, let
\begin{equation}
  Y_p
  \defeq
  \overline{
    h_\omega^{\frac12(\frac1p-\frac1r)}
    Y_r
    h_\omega^{\frac12(\frac1p-\frac1r)}}
  \in L^p(M)_+
  \label{eq:anchored-entropy-path}
\end{equation}
be the compatible path of
Lemma~\ref{lem:compatible-weighted-scale}, and set
\begin{equation}
  Z(p)\defeq\tr(Y_p^p),
  \qquad
  G(p)\defeq\frac1p\log Z(p)=\log\|Y_p\|_p,
  \qquad
  h_{\eta_p}\defeq\frac{Y_p^p}{Z(p)}.
  \label{eq:anchored-entropy-data}
\end{equation}
Then $Z(p)\in(0,\infty)$ and $\eta_p$ is a normal state.  Moreover:
\begin{enumerate}[label=\textup{(\roman*)}]
\item For every $0<p'<p\leq r$,
\begin{align}
  Q_{p'/p,p'}(\eta_p\|\omega)
  &=Z(p')Z(p)^{-p'/p},
  \label{eq:anchored-exact-quotient}\\
  \frac{G(p')-G(p)}{p'-p}
  &=\frac{D_{p'/p,p'}(\eta_p\Vert\omega)}{pp'}.
  \label{eq:anchored-exact-difference-quotient}
\end{align}

\item For every $0<p\leq r$, the left derivative exists in the
extended-real sense and satisfies
\begin{equation}
  G'_-(p)
  =\frac{D(\eta_p\Vert\omega)}{p^2}.
  \label{eq:anchored-entropy-left-derivative}
\end{equation}
For $p<r$ both sides are finite.  At the anchoring endpoint $p=r$
they may equal $+\infty$.

\item The map
\begin{equation}
  p\longmapsto D(\eta_p\Vert\omega)
  \label{eq:anchored-entropy-profile}
\end{equation}
is nondecreasing and Borel measurable as an
$[0,+\infty]$-valued function on $(0,r]$.

\item The function $G$ is absolutely continuous on every compact interval
$[p_0,p_1]\subset(0,r]$, and
\begin{equation}
  G(p_1)-G(p_0)
  =
  \int_{p_0}^{p_1}
  \frac{D(\eta_p\Vert\omega)}{p^2}\,\dd p,
  \qquad 0<p_0<p_1\leq r.
  \label{eq:anchored-entropy-integral}
\end{equation}
In particular,
\begin{equation}
  \log\frac{\|Y_r\|_r}{\|Y_{p_0}\|_{p_0}}
  =
  \int_{p_0}^{r}
  \frac{D(\eta_p\Vert\omega)}{p^2}\,\dd p.
  \label{eq:anchored-entropy-integral-to-anchor}
\end{equation}
These integrals are finite.  If
$D(\eta_r\Vert\omega)=+\infty$, its value at the single endpoint $r$
does not affect the Lebesgue integral.
\end{enumerate}
\end{proposition}

\begin{proof}
Lemma~\ref{lem:compatible-weighted-scale} gives
$0\ne Y_p\in L^p(M)_+$.  Hence
$Z(p)\in(0,\infty)$ and $Z(p)^{-1}Y_p^p$ is the Haagerup density of a
normal state.

Fix $0<p'<p\leq r$ and put
\[
  a\defeq\frac{p'}p\in(0,1).
\]
Functional calculus gives
\begin{align}
  \overline{
  h_\omega^{\frac{1-a}{2p'}}
  h_{\eta_p}^{a/p'}
  h_\omega^{\frac{1-a}{2p'}}}
  =
  Z(p)^{-1/p}
  \overline{
    h_\omega^{\frac12(\frac1{p'}-\frac1p)}
    Y_p
    h_\omega^{\frac12(\frac1{p'}-\frac1p)}}.
\end{align}
Compatibility of the anchored path identifies the closed product on the
right with $Y_{p'}$.  Taking the $p'$-th moment therefore gives
\eqref{eq:anchored-exact-quotient}.  Since $\eta_p$ and $\omega$ are
states,
\begin{align}
  D_{p'/p,p'}(\eta_p\Vert\omega)
  &=
  \frac{1}{p'/p-1}
  \log\!\left(Z(p')Z(p)^{-p'/p}\right)\\
  &=
  pp'\frac{G(p')-G(p)}{p'-p},
\end{align}
which proves \eqref{eq:anchored-exact-difference-quotient}.

We now let $p'\uparrow p$.  For $p'$ sufficiently close to $p$, so that
$p/2<p'<p$, monotonicity in the second parameter at lower order
\cite[Theorem~1 (x)]{Kato:2023hlj} gives
\begin{equation}
  D_{p'/p,p/2}(\eta_p\Vert\omega)
  \leq
  D_{p'/p,p'}(\eta_p\Vert\omega)
  \leq
  D_{p'/p,p}(\eta_p\Vert\omega).
  \label{eq:anchored-varying-z-squeeze}
\end{equation}
For the fixed second parameters $p/2$ and $p$, the lower-order endpoint
theorem \cite[Theorem~6.7]{Hiai:2024qve} yields
\begin{equation}
  \lim_{p'\uparrow p}D_{p'/p,p/2}(\eta_p\Vert\omega)
  =
  \lim_{p'\uparrow p}D_{p'/p,p}(\eta_p\Vert\omega)
  =
  D(\eta_p\Vert\omega)
\end{equation}
in $[0,+\infty]$.  Squeezing in
\eqref{eq:anchored-exact-difference-quotient} proves
\eqref{eq:anchored-entropy-left-derivative}.

Define
\begin{equation}
  B(u)\defeq G(1/u)=\log\|Y_{1/u}\|_{1/u},
  \qquad u\geq\frac1r.
\end{equation}
By Lemma~\ref{lem:compatible-weighted-scale}, $B$ is finite and convex on
$(1/r,\infty)$ and continuous at $1/r$.  The one-sided chain rule gives
\begin{equation}
  B'_+(1/p)
  =
  -p^2G'_-(p)
  =
  -D(\eta_p\Vert\omega).
  \label{eq:anchored-B-derivative}
\end{equation}
If $p<r$, then $1/p$ is an interior point of the domain of the finite
convex function $B$, so its one-sided derivative, and hence
$D(\eta_p\Vert\omega)$, is finite.

Since $B'_+$ is nondecreasing in $u$, reversing the reciprocal variable
in \eqref{eq:anchored-B-derivative} shows that
$p\mapsto D(\eta_p\Vert\omega)$ is nondecreasing.  It is therefore Borel
measurable.

Finally, Lemma~\ref{lem:convex-endpoint-ftc} shows that $B$ is absolutely
continuous on every compact interval $[1/r,b]$.  Since $p\mapsto1/p$ is
bi-Lipschitz on compact subintervals of $(0,r]$, $G(p)=B(1/p)$ is
absolutely continuous on $[p_0,p_1]$.  Its ordinary derivative agrees
almost everywhere with the left derivative in
\eqref{eq:anchored-entropy-left-derivative}.  The fundamental theorem of
calculus proves \eqref{eq:anchored-entropy-integral}; taking $p_1=r$
gives \eqref{eq:anchored-entropy-integral-to-anchor}.  Finiteness follows
from the finiteness of the endpoint values of $G$.
\end{proof}

\subsection{Fixed-ray theorem as a specialization of the anchored identity}
\label{subsec:fixed-ray-proof}

Proposition~\ref{prop:anchored-entropy-identity} is 
independent of any specific values of
$\alpha$, $z$ beyond positivity, and of the fixed-ray construction.  We now show that
Theorem~\ref{thm:fixed-ray-integral} is obtained by choosing a
branch-dependent anchor in that proposition. The proof has two parts. 
We first prove the factor-existence and
integral-representation assertions of Theorem~\ref{thm:fixed-ray-integral}.  We then prove the
monotonicity and one-sided derivative assertions of Corollary~\ref{cor:raywise-monotonicity}.

Put $e=s(\omega)$, pass to the corner $eMe$, and, for the duration of the
proof, relabel the corner and the restricted states as $M$, $\psi$, and
$\omega$.  Thus $\omega$ is faithful.  The states constructed in the
corner extend canonically to the original algebra, and the corresponding
relative entropies are unchanged.

Recall that $c=z/\alpha$.  Choose the anchor
\begin{equation}
  (r,Y_r)\defeq
  \begin{cases}
    \bigl(c,h_\psi^{1/c}\bigr),
      &0<\alpha<1,\\[0.5ex]
    \bigl(z,x_{\alpha,z}(\psi\|\omega)\bigr),
      &\alpha>1.
  \end{cases}
  \label{eq:fixed-ray-anchor}
\end{equation}
The lower anchor is available by the support-inclusion hypothesis.  In
the upper branch, the factor exists and belongs to $L^z(M)_+$ by the
finiteness hypothesis.

Let $(Y_p)_{0<p\leq r}$ be the associated anchored path, and use the
notation $Z(p)$, $G(p)$, and $\eta_p$ of Proposition~\ref{prop:anchored-entropy-identity}.  In the lower branch, for $\beta\in[\alpha,1]$,
\[
    Y_{c\beta} = J_{c\beta,c}^{\omega}(h_\psi^{1/c}) = h_\omega^{\frac{1-\beta}{2c\beta}} h_\psi^{1/c}h_\omega^{\frac{1-\beta}{2c\beta}} = x_\beta^{(c)}(\psi\|\omega).
\]
In the upper branch, 
the defining factorization \eqref{eq:upper-fixed-ray-factor} for the anchor $Y_z=x_{\alpha,z}(\psi\|\omega)$ is
\begin{equation}
  h_\psi^{1/c}
  =
  \overline{
  h_\omega^{\frac1{2c}-\frac1{2z}}
  Y_z
  h_\omega^{\frac1{2c}-\frac1{2z}}}
  =
  J_{c,z}^\omega(Y_z).
  \label{eq:upper-fixed-ray-anchor-factorization}
\end{equation}
For $\beta\in[1,\alpha]$, we have
\begin{equation}
    h_\psi^{1/c}
    =
    J_{c,z}^\omega(Y_z)
    = 
    J_{c,c\beta}^\omega J_{c\beta,z}^\omega(Y_z)
    = 
    J_{c,c\beta}^\omega(Y_{c\beta})
    =
    h_\omega^{\frac{\beta-1}{2c\beta}}Y_{c\beta}h_\omega^{\frac{\beta-1}{2c\beta}}.
\end{equation}
By uniqueness, $Y_{c\beta}=x_\beta^{(c)}$ also in the upper branch.

Taking powers, traces and normalized densities then gives
\begin{equation}
  Z(c\beta)
  =
  Q_{\beta,c\beta}(\psi\|\omega),
  \qquad
  \eta_{c\beta}
  =
  \psi_\beta^{(c)},
  \qquad \beta\in I_\alpha.
  \label{eq:anchored-fixed-ray-identification}
\end{equation}
Thus all the intermediate fixed-ray factors have finite nonzero moments.

This proves the existence and finite-nonzero-moment assertions of
Theorem~\ref{thm:fixed-ray-integral} and identifies 
the normalized powers with the fixed-ray
escort states.  It remains to prove the integral identity~\eqref{eq:main-fixed-ray-integral}.

At the two distinguished endpoints,
\begin{equation}
  Z(c)=\tr(h_\psi)=1,
  \qquad
  Z(z)=Q_{\alpha,z}(\psi\|\omega),
  \qquad
  G(c)=0.
  \label{eq:fixed-ray-anchor-endpoints}
\end{equation}

Applying \eqref{eq:anchored-entropy-integral} between $p=c$ and $p=z$
gives
\begin{align}
  \frac1z\log Q_{\alpha,z}(\psi\|\omega)
  &=
  G(z)-G(c)\notag\\
  &=
  \int_c^z
  \frac{D(\eta_p\Vert\omega)}{p^2}\,\dd p\notag\\
  &=
  \frac1c
  \int_1^\alpha
  \frac{D(\psi_\beta^{(c)}\Vert\omega)}{\beta^2}\,\dd\beta.
  \label{eq:fixed-ray-from-anchored-identity}
\end{align}
Here the integrals are oriented; when $0<\alpha<1$, this identity is
obtained by applying
\eqref{eq:anchored-entropy-integral} on $[z,c]$ and reversing the signs.
Since $z=c\alpha$, multiplication by $z/(\alpha-1)$ yields
\begin{equation}
  D_{\alpha,z}(\psi\Vert\omega)
  =
  \frac{\alpha}{\alpha-1}
  \int_1^\alpha
  \frac{D(\psi_\beta^{(c)}\Vert\omega)}{\beta^2}\,\dd\beta.
\end{equation}
This is \eqref{eq:main-fixed-ray-integral}; the formulation using
$\mu_\alpha$ is the same identity with the orientation removed.  Thus the
fixed-ray representation is a specialization of the general anchored
entropy identity.

This completes the proof of Theorem~\ref{thm:fixed-ray-integral}; 
Corollary~\ref{cor:srd-integral} follows by
setting \(c=1\).  We now prove Corollary~\ref{cor:raywise-monotonicity}, 
beginning with the monotonicity of \(R_c\).  On every anchored
segment\footnote{Here an \emph{anchored segment} means an interval of the fixed-ray
parameter on which the factors are obtained from a single anchored
path.}, Proposition~\ref{prop:anchored-entropy-identity} shows that
$p\mapsto D(\eta_p\Vert\omega)$ is nondecreasing.  Since $p=c\beta$,
\eqref{eq:anchored-fixed-ray-identification} implies that
$\beta\mapsto R_c(\beta)$ is nondecreasing on that segment.

Let us now prove \eqref{eq:raywise-refined-identity}. These anchored 
segments cover all of $\mathcal J_c$.  Indeed, if
$0<\beta_1<\beta_2\leq1$, use the lower path anchored at $p=c$.  If
$1\leq\beta_1<\beta_2$ and $\beta_2\in\mathcal J_c$, use the upper path
anchored at $p=c\beta_2$.  If $\beta_1<1<\beta_2$, combine the preceding
two comparisons through $R_c(1)$.  Hence $R_c$ is nondecreasing on
$\mathcal J_c$. As a monotone extended-real-valued function 
on the interval \(\mathcal J_c\), \(R_c\) is Borel measurable.

On any such anchored segment,
\begin{equation}
  F_c(\beta)
  =
  \frac1\beta\log Z(c\beta)
  =
  cG(c\beta).
  \label{eq:fixed-ray-potential-from-anchor}
\end{equation}
Therefore \eqref{eq:anchored-entropy-left-derivative} and the one-sided
chain rule give
\begin{align}
  \beta^2\frac{\dd^-}{\dd\beta}F_c(\beta)
  &=
  \beta^2c^2G'_-(c\beta)\notag\\
  &=
  D(\eta_{c\beta}\Vert\omega)
  =
  R_c(\beta).
\end{align}
If \(\beta>1\) is an interior point of \(\mathcal J_c\), choose
\(\gamma\in\mathcal J_c\) with \(\gamma>\beta\).  Applying
Proposition~\ref{prop:anchored-entropy-identity} to the upper path anchored at \(r=c\gamma\) places
\(p=c\beta\) strictly below the anchor.  Hence
\(R_c(\beta)=D(\eta_{c\beta}\Vert\omega)\) is finite, and the derivative
identity applies there as an interior identity.
Points below one belong to the lower path.  At
$\beta=1$, the lower path gives the asserted endpoint left derivative,
with the extended-real interpretation when
$D(\psi\Vert\omega)=+\infty$.  

If \(M\) is finite-dimensional, then after passing to the support
corner \(h_\omega\) is invertible and
\[
  x_\beta^{(c)}
  =
  h_\omega^{(1-\beta)/(2c\beta)}
  h_\psi^{1/c}
  h_\omega^{(1-\beta)/(2c\beta)}
\]
is a smooth positive-matrix family of constant rank. Smooth
finite-dimensional functional calculus therefore shows that
\(F_c\) and \(R_c\) are smooth on the interior of \(J_c\), so the
left derivative there agrees with the ordinary derivative.
This completes the proof of
Corollary~\ref{cor:raywise-monotonicity}.

\subsection{Proof of the general lower-order boundary formula}
\label{subsec:lower-order-proof}

It remains to prove Proposition~\ref{prop:general-fixed-ray-boundary}.
We first prove the bound \eqref{eq:boundary-cost-bound}, 
then treat the case \(q_c>0\), proving
the escort identification and the boundary 
formula \eqref{eq:general-fixed-ray-boundary}, and finally
handle the case \(q_c=0\).

Let $e=s(\omega)$ and use the canonical corner identification.  Since
$x\mapsto{exe}$ is a positive contraction on $L^c(M)$,
$y_c\in L^c(eMe)_+$ and $q_c<\infty$.  Faithfulness of the canonical
functional on $L^1(eMe)_+$ also gives $q_c=0$ only if $y_c=0$.
Moreover, contractivity of the compression gives
\[
  q_c
  = \|y_c\|_c^c
  = \bigl\|e h_\psi^{1/c}e\bigr\|_c^c
  \leq \bigl\|h_\psi^{1/c}\bigr\|_c^c
  = \tr(h_\psi)
  = 1.
\]
This proves \eqref{eq:boundary-cost-bound}, i.e.\ \(0\leq q_c\leq1\).


Since
$h_\omega$ is supported on $e$, for $0<\beta<1$ we may insert $e$ on both
sides of $h_\psi^{1/c}$ in the closed product:
\begin{align}
  x_\beta^{(c)}(\psi\|\omega)
  &=\overline{
  h_\omega^{\frac{1-\beta}{2c\beta}}
  y_c
  h_\omega^{\frac{1-\beta}{2c\beta}}}.
  \label{eq:general-boundary-factor-compression}
\end{align}
If $q_c>0$, then
$y_c=q_c^{1/c}h_{\widehat\psi_c}^{1/c}$, and hence
\begin{equation}
  x_\beta^{(c)}(\psi\|\omega)
  =q_c^{1/c}
   x_\beta^{(c)}(\widehat\psi_c\|\omega_e).
  \label{eq:general-boundary-factor-scaling}
\end{equation}
Taking the $c\beta$-th moment gives
\begin{equation}
  Q_{\beta,c\beta}(\psi\|\omega)
  =q_c^\beta
   Q_{\beta,c\beta}(\widehat\psi_c\|\omega_e).
  \label{eq:general-boundary-moment-scaling}
\end{equation}
The factor \(q_c^\beta\) in \eqref{eq:general-boundary-moment-scaling} cancels upon normalization.
Consequently, \eqref{eq:general-boundary-factor-scaling} and \eqref{eq:general-boundary-moment-scaling} prove the escort-identification assertion
of Proposition~\ref{prop:general-fixed-ray-boundary}.  

For \(q_c>0\) at $\beta=\alpha$,
\begin{equation}
  D_{\alpha,c\alpha}(\psi\Vert\omega)
  =D_{\alpha,c\alpha}(\widehat\psi_c\Vert\omega_e)
   -\frac{\alpha}{1-\alpha}\log q_c.
\end{equation}
The state $\widehat\psi_c$ is supported in $e$, so applying
Theorem~\ref{thm:fixed-ray-integral} to the corner proves
\eqref{eq:general-fixed-ray-boundary} when \(q_c>0\).  

If $q_c=0$, then $y_c=0$, and
\eqref{eq:general-boundary-factor-compression} gives
\(Q_{\alpha,c\alpha}(\psi\Vert\omega)=0\), hence
\(D_{\alpha,c\alpha}(\psi\Vert\omega)=+\infty\).
This proves the remaining case and completes the proof of Proposition~\ref{prop:general-fixed-ray-boundary}.
Corollary~\ref{cor:srd-support-boundary} is its specialization to \(c=1\).

\section{Discussion}
\label{sec:discussion}

In this paper we have shown that, for supported pairs of states $\psi,\omega$ such that $s(\psi)\leq s(\omega)$, the {\az} divergence $D_{\alpha,z}(\psi\|\omega)$ is a probability average of ordinary relative entropies along a canonical fixed-ray escort trajectory. For lower orders, i.e., $\alpha\in(0,1)$, without support inclusion, the corresponding compressed average is supplemented by the explicit $L^c$ support-boundary term.
This representation holds both for density matrices in finite-dimensional systems and for normal states on arbitrary von Neumann algebras, and needs no DPI restriction.
We have also discussed operational interpretations of our representation through independent theorems: generic DPI parameters give testing
converses, the sandwiched line gives exact strong-converse and
thermal-operation reliability formulas, and a restricted pair conversion
problem uses the full two-parameter family. 

In a companion work~\cite{Kibe:2026hepth}, we apply the integral representation to null shape deformations in quantum field theory (QFT) and obtain, under the regulated-QFT assumptions stated there, quasi-local expressions for \(\alpha\)--\(z\) R\'enyi quantum null energy condition quantities as escort averages of ordinary relative-entropy variations. 
Two further consequences of the fixed-ray construction will be developed elsewhere. Under suitable regularity assumptions, differentiating the associated escort path a second time leads to a modular susceptibility, a quadratic fluctuation quantity interpolating between relative-entropy variance and Bogoliubov-Kubo-Mori covariance. The escort composition law also suggests invariance of quantum sufficiency along an admissible fixed ray and could lead to simultaneous recovery estimates for an associated order-labelled escort experiment. 



Another interesting direction for further research is to compare the escort representation with other
integral descriptions of R\'enyi divergences.  Hirche and Tomamichel 
construct quantum $f$-divergences by integrating quantum hockey-stick
divergences and show that the regularized R\'enyi quantities recover the
Petz R\'enyi divergence below one and the sandwiched R\'enyi divergence
above one \cite{Hirche:2023caq}.  Their construction resolves the
divergence into binary-testing quantities, whereas the representation
developed here resolves it into ordinary relative entropies evaluated along
a canonical path of states.  There is also a classical precedent for the
latter viewpoint: van Erven and Harremo\"es organize classical R\'enyi
divergence using the tilted distributions
which appear naturally in variational identities involving
Kullback--Leibler divergences \cite{Erven:2014aa}.  It would be
interesting to determine whether the hockey-stick and escort-path
representations are related by a direct transform, and to what extent the
sandwiched escort family provides the intrinsic noncommutative analogue of
the classical tilted-distribution path.

\begingroup
\renewcommand{\addcontentsline}[3]{}
\section*{Acknowledgments}

T.K. is supported by a Simons Foundation fellowship through the Targeted Grant to Instituto Balseiro. The work of P.R. has been supported by the Polish National Science Centre through Sonata grant (2022/47/D/ST2/02058). 
During the preparation of this manuscript, the authors used Anthropic Claude and OpenAI ChatGPT as editorial aids in drafting and revising portions of the text. Every mathematical statement, proof, and citation in the final manuscript was independently checked by the authors, who take full responsibility for its content.

\endgroup


\providecommand{\href}[2]{#2}\begingroup\raggedright\begin{thebibliography}{10}

\bibitem{Araki:1976zv}
H.~Araki, \emph{{Relative Entropy of States of von Neumann Algebras}},
  \href{https://doi.org/10.2977/prims/1195191148}{\emph{Publ. Res. Inst. Math.
  Sci.} {\bfseries 11} (1976) 809}.

\bibitem{Muller-Lennert:2013liu}
M.~M{\"u}ller-Lennert, F.~Dupuis, O.~Szehr, S.~Fehr and M.~Tomamichel,
  \emph{{On quantum R{\'e}nyi entropies: A new generalization and some
  properties}}, \href{https://doi.org/10.1063/1.4838856}{\emph{J. Math. Phys.}
  {\bfseries 54} (2013) 122203}
  [\href{https://arxiv.org/abs/1306.3142}{{\ttfamily 1306.3142}}].

\bibitem{Wilde:2013bdg}
M.M.~Wilde, A.~Winter and D.~Yang, \emph{{Strong Converse for the Classical
  Capacity of Entanglement-Breaking and Hadamard Channels via a Sandwiched
  Renyi Relative Entropy}},
  \href{https://doi.org/10.1007/s00220-014-2122-x}{\emph{Commun. Math. Phys.}
  {\bfseries 331} (2014) 593}
  [\href{https://arxiv.org/abs/1306.1586}{{\ttfamily 1306.1586}}].

\bibitem{Audenaert:2015npv}
K.M.R.~Audenaert and N.~Datta, \emph{{alpha-z-relative Renyi entropies}},
  \href{https://doi.org/10.1063/1.4906367}{\emph{J. Math. Phys.} {\bfseries 56}
  (2015) 022202} [\href{https://arxiv.org/abs/1310.7178}{{\ttfamily
  1310.7178}}].

\bibitem{Kato:2023aro}
S.~Kato and Y.~Ueda, \emph{{A Remark on Non-commutative $L^p$-Spaces}},
  \href{https://doi.org/10.4064/sm230724-11-10}{\emph{Studia Math.} {\bfseries
  275} (2024) 235} [\href{https://arxiv.org/abs/2307.01790}{{\ttfamily
  2307.01790}}].

\bibitem{Kato:2023hlj}
S.~Kato, \emph{{On {\ensuremath{\alpha}}-z-R{\'e}nyi divergence in the von
  Neumann algebra setting}}, \href{https://doi.org/10.1063/5.0186552}{\emph{J.
  Math. Phys.} {\bfseries 65} (2024) 042202}
  [\href{https://arxiv.org/abs/2311.01748}{{\ttfamily 2311.01748}}].

\bibitem{Hiai:2024qve}
F.~Hiai and A.~Jen{\v{c}}ov{\'a}, \emph{{$\alpha $-z-R{\'e}nyi Divergences in
  von Neumann Algebras: Data Processing Inequality, Reversibility, and
  Monotonicity Properties in $\alpha ,z$}},
  \href{https://doi.org/10.1007/s00220-024-05124-1}{\emph{Commun. Math. Phys.}
  {\bfseries 405} (2024) 271}
  [\href{https://arxiv.org/abs/2404.07617}{{\ttfamily 2404.07617}}].

\bibitem{Lashkari:2018nsl}
N.~Lashkari, \emph{{Constraining Quantum Fields using Modular Theory}},
  \href{https://doi.org/10.1007/JHEP01(2019)059}{\emph{JHEP} {\bfseries 01}
  (2019) 059} [\href{https://arxiv.org/abs/1810.09306}{{\ttfamily
  1810.09306}}].

\bibitem{Moosa:2020jwt}
M.~Moosa, P.~Rath and V.P.~Su, \emph{{A R\'enyi quantum null energy condition:
  proof for free field theories}},
  \href{https://doi.org/10.1007/JHEP01(2021)064}{\emph{JHEP} {\bfseries 01}
  (2021) 064} [\href{https://arxiv.org/abs/2007.15025}{{\ttfamily
  2007.15025}}].

\bibitem{Roy:2022yzm}
P.~Roy, \emph{{Proof of the R\'enyi quantum null energy condition for free
  fermions}}, \href{https://doi.org/10.1103/PhysRevD.108.045010}{\emph{Phys.
  Rev. D} {\bfseries 108} (2023) 045010}
  [\href{https://arxiv.org/abs/2212.02331}{{\ttfamily 2212.02331}}].

\bibitem{Kibe:2026wsg}
T.~Kibe and P.~Roy, \emph{{A general proof of integer R{\'e}nyi QNEC}},
  \href{https://arxiv.org/abs/2605.15272}{{\ttfamily 2605.15272}}.

\bibitem{Kibe:2026bcn}
T.~Kibe and P.~Roy, \emph{{No off-diagonal quantum focusing for R{\'e}nyi
  divergences}},  \href{https://arxiv.org/abs/2607.07799}{{\ttfamily
  2607.07799}}.

\bibitem{Bao:2019aol}
N.~Bao, M.~Moosa and I.~Shehzad, \emph{{The holographic dual of R\'enyi
  relative entropy}},
  \href{https://doi.org/10.1007/JHEP08(2019)099}{\emph{JHEP} {\bfseries 08}
  (2019) 099} [\href{https://arxiv.org/abs/1904.08433}{{\ttfamily
  1904.08433}}].

\bibitem{Hirche:2023caq}
C.~Hirche and M.~Tomamichel, \emph{{Quantum R{\'e}nyi and f-Divergences from
  Integral Representations}},
  \href{https://doi.org/10.1007/s00220-024-05087-3}{\emph{Commun. Math. Phys.}
  {\bfseries 405} (2024) 208}
  [\href{https://arxiv.org/abs/2306.12343}{{\ttfamily 2306.12343}}].

\bibitem{Liu:2025npj}
P.-C.~Liu, C.~Hirche and H.-C.~Cheng, \emph{{Layer Cake Representations for
  Quantum Divergences}},  \href{https://arxiv.org/abs/2507.07065}{{\ttfamily
  2507.07065}}.

\bibitem{Gu:2019jfc}
J.~Gu, Z.~Yin and H.~Zhang, \emph{{Interpolation of quasi noncommutative
  $L_p$-spaces}},  \href{https://arxiv.org/abs/1905.08491}{{\ttfamily
  1905.08491}}.

\bibitem{Bhatia:1997}
R.~Bhatia, \emph{Matrix Analysis}, vol.~169, Springer, Berlin (1997).

\bibitem{Terp:1981lp}
M.~Terp, \emph{{$L^p$} spaces associated with von Neumann algebras},
  Math.~Inst.~Copenhagen University (1981).

\bibitem{KOSAKI198429}
H.~Kosaki, \emph{{Applications of the Complex Interpolation Method to a von
  Neumann Algebra: Non-commutative $L^p$-Spaces}},
  \href{https://doi.org/10.1016/0022-1236(84)90025-9}{\emph{J. Funct. Anal.}
  {\bfseries 56} (1984) 29}.

\bibitem{Berta:2016vnw}
M.~Berta, V.B.~Scholz and M.~Tomamichel, \emph{{R\'enyi Divergences as Weighted
  Non-commutative Vector-Valued $L_p$ -Spaces}},
  \href{https://doi.org/10.1007/s00023-018-0670-x}{\emph{Annales Henri
  Poincare} {\bfseries 19} (2018) 1843}
  [\href{https://arxiv.org/abs/1608.05317}{{\ttfamily 1608.05317}}].

\bibitem{Jencova:2016tqz}
A.~Jen{\v{c}}ov{\'a}, \emph{{R{\'e}nyi Relative Entropies and Noncommutative
  $L_p$-Spaces}},
  \href{https://doi.org/10.1007/s00023-018-0683-5}{\emph{Annales Henri
  Poincar{\'e}} {\bfseries 19} (2018) 2513}
  [\href{https://arxiv.org/abs/1609.08462}{{\ttfamily 1609.08462}}].

\bibitem{Jencova:2017txf}
A.~Jen\v{c}ov\'a, \emph{{R\'enyi relative entropies and noncommutative
  $L_p$-spaces II}},
  \href{https://doi.org/10.1007/s00023-021-01074-9}{\emph{Annales Henri
  Poincare} {\bfseries 22} (2021) 3235}
  [\href{https://arxiv.org/abs/1707.00047}{{\ttfamily 1707.00047}}].

\bibitem{Zhang:2018roy}
H.~Zhang, \emph{{From Wigner-Yanase-Dyson conjecture to Carlen-Frank-Lieb
  conjecture}}, \href{https://doi.org/10.1016/j.aim.2020.107053}{\emph{Adv.
  Math.} {\bfseries 365} (2020) 107053}
  [\href{https://arxiv.org/abs/1811.01205}{{\ttfamily 1811.01205}}].

\bibitem{Watanabe:2026pzw}
K.~Watanabe, B.~Regula, M.~Tomamichel and R.~Takagi, \emph{{Reliability of
  asymptotic work extraction}},
  \href{https://arxiv.org/abs/2606.06318}{{\ttfamily 2606.06318}}.

\bibitem{Hiai1981SufficiencyKC}
F.~Hiai, M.~Ohya and M.~Tsukada, \emph{Sufficiency, kms condition and relative
  entropy in von neumann algebras.}, {\emph{Pacific Journal of Mathematics}
  {\bfseries 96} (1981) 99}.

\bibitem{Holevo:2004ots}
A.S.~Holevo and M.E.~Shirokov, \emph{{Continuous ensembles and the
  $\chi$-capacity of infinite-dimensional channels}}, {\emph{Theor. Probab.
  Appl.} {\bfseries 50} (2005) 86}
  [\href{https://arxiv.org/abs/quant-ph/0408176}{{\ttfamily
  quant-ph/0408176}}].

\bibitem{Mosonyi:2014wsq}
M.~Mosonyi and T.~Ogawa, \emph{{Quantum Hypothesis Testing and the Operational
  Interpretation of the Quantum R{\'e}nyi Relative Entropies}},
  \href{https://doi.org/10.1007/s00220-014-2248-x}{\emph{Commun. Math. Phys.}
  {\bfseries 334} (2015) 1617}
  [\href{https://arxiv.org/abs/1309.3228}{{\ttfamily 1309.3228}}].

\bibitem{Hiai:2021slp}
F.~Hiai and M.~Mosonyi, \emph{{Quantum R{\'e}nyi Divergences and the Strong
  Converse Exponent of State Discrimination in Operator Algebras}},
  \href{https://doi.org/10.1007/s00023-022-01250-5}{\emph{Annales Henri
  Poincare} {\bfseries 24} (2023) 1681}
  [\href{https://arxiv.org/abs/2110.07320}{{\ttfamily 2110.07320}}].

\bibitem{Junge:2025kuz}
M.~Junge and N.~LaRacuente, \emph{{Strong Converse Rate for Asymptotic
  Hypothesis Testing in Type III}},
  \href{https://arxiv.org/abs/2507.07989}{{\ttfamily 2507.07989}}.

\bibitem{Verhagen:2025anx}
F.~Verhagen, M.~Tomamichel and E.~Haapasalo, \emph{{Conditions for Large-Sample
  Majorization of Pairs of Flat States in Terms of $\alpha $-z Relative
  Entropies}}, \href{https://doi.org/10.1007/s00220-026-05675-5}{\emph{Commun.
  Math. Phys.} {\bfseries 407} (2026) 157}
  [\href{https://arxiv.org/abs/2507.07520}{{\ttfamily 2507.07520}}].

\bibitem{Araki:1977zsq}
H.~Araki, \emph{{Relative Entropy for States of von Neumann Algebras II}},
  \href{https://doi.org/10.2977/prims/1195190105}{\emph{Publ. Res. Inst. Math.
  Sci. Kyoto} {\bfseries 13} (1977) 173}.

\bibitem{Junge:2003nbr}
M.~Junge and Q.~Xu, \emph{{Noncommutative Burkholder/Rosenthal inequalities}},
  \href{https://doi.org/10.1214/aop/1048516542}{\emph{Annals Probab.}
  {\bfseries 31} (2003) 948}.

\bibitem{Fack:1986ezo}
T.~Fack and H.~Kosaki, \emph{Generalized s-numbers of $\tau$-measurable
  operators}, \href{https://doi.org/10.2140/pjm.1986.123.269}{\emph{Pacific
  Journal of Mathematics} {\bfseries 123} (1986) 269}.

\bibitem{Kibe:2026hepth}
T.~Kibe and P.~Roy, \emph{To appear soon},  2026.

\bibitem{Erven:2014aa}
T.~van Erven and P.~Harremo{\"e}s, \emph{{R{\'e}nyi Divergence and
  Kullback--Leibler Divergence}},
  \href{https://doi.org/10.1109/TIT.2014.2320500}{\emph{IEEE Trans. Inf.
  Theory} {\bfseries 60} (2014) 3797}
  [\href{https://arxiv.org/abs/1206.2459}{{\ttfamily 1206.2459}}].

\end{thebibliography}\endgroup


\end{document}
